% Springer-compatible LaTeX template for APJM / MIR submission
% Compiled from Markdown via pandoc.
% Document class: `article` (Springer Nature's sn-jnl is not bundled; this template
% follows Springer's general manuscript guidelines and is compatible with both
% Springer Nature LaTeX template and Word-track submission portals).

\documentclass[12pt,a4paper]{article}

% --- Geometry & spacing ---
\usepackage[margin=2.5cm]{geometry}
\usepackage{setspace}
\onehalfspacing

% --- Encoding & fonts ---
% Conditional setup: under pdflatex use legacy inputenc; under xelatex/lualatex
% use fontspec for native Unicode support (needed for β, −, × symbols).
\usepackage{iftex}
\ifPDFTeX
  \usepackage[utf8]{inputenc}
  \usepackage[T1]{fontenc}
  \usepackage{lmodern}
\else
  \usepackage{fontspec}
  \setmainfont{Liberation Serif}[
    Scale=1.0,
    BoldFont={Liberation Serif Bold},
    ItalicFont={Liberation Serif Italic},
    BoldItalicFont={Liberation Serif Bold Italic}
  ]
  % Liberation Mono carries the full Vietnamese diacritic set; DejaVu Sans
  % Mono does not (renders Ấ/ầ/ế/ố as missing characters in code blocks).
  \setmonofont{Liberation Mono}[Scale=0.85]
  % Map common Unicode math/typographic symbols that may appear in markdown
  % body text. Liberation Serif lacks them, so route via amssymb/math mode.
  % (We deliberately avoid unicode-math, which conflicts with the amssymb
  % package that pandoc auto-loads in the preamble.)
  \usepackage{newunicodechar}
  \newunicodechar{∈}{\ensuremath{\in}}
  \newunicodechar{∉}{\ensuremath{\notin}}
  \newunicodechar{≤}{\ensuremath{\leq}}
  \newunicodechar{≥}{\ensuremath{\geq}}
  \newunicodechar{≠}{\ensuremath{\neq}}
  \newunicodechar{≈}{\ensuremath{\approx}}
  \newunicodechar{∞}{\ensuremath{\infty}}
  \newunicodechar{∑}{\ensuremath{\sum}}
  \newunicodechar{∏}{\ensuremath{\prod}}
  \newunicodechar{∗}{\ensuremath{\ast}}
  \newunicodechar{·}{\ensuremath{\cdot}}
  \newunicodechar{−}{\ensuremath{-}}
  \newunicodechar{×}{\ensuremath{\times}}
  \newunicodechar{÷}{\ensuremath{\div}}
  \newunicodechar{±}{\ensuremath{\pm}}
  % Typographic checkmarks/crosses (used as bullet markers in markdown).
  \newunicodechar{✓}{\ensuremath{\checkmark}}
  \newunicodechar{✗}{\ensuremath{\times}}
  \newunicodechar{⚠}{\textbf{!}}
  \newunicodechar{→}{\ensuremath{\rightarrow}}
  \newunicodechar{←}{\ensuremath{\leftarrow}}
  \newunicodechar{⇒}{\ensuremath{\Rightarrow}}
  % Allow URL line breaks at slashes (fix Overfull \hbox in long URLs)
  \usepackage{xurl}
\fi

% --- Math ---
\usepackage{amsmath,amssymb,amsfonts,mathtools}
\usepackage{bm}

% --- Tables ---
\usepackage{booktabs}
\usepackage{longtable}
\usepackage{multirow}
\usepackage{array}
\usepackage{tabularx}

% --- Figures ---
\usepackage{graphicx}
\usepackage{float}
\usepackage{caption}
\captionsetup{font=small,labelfont=bf}

% --- Misc text ---
\usepackage{microtype}
\usepackage{enumitem}
\setlist{nosep,leftmargin=1.5em}
\usepackage{textcomp}
\usepackage{xcolor}
\usepackage{ulem}\normalem

% --- Code block support (pandoc Shaded / Highlighting environments) ---
% Required when source markdown contains fenced code blocks (```...```).
\usepackage{fancyvrb}
\usepackage{framed}
\definecolor{shadecolor}{RGB}{248,248,248}
\newenvironment{Shaded}{\begin{snugshade}}{\end{snugshade}}
% Highlighting environment for syntax-highlighted code
\providecommand{\BuiltInTok}[1]{#1}
\providecommand{\CommentTok}[1]{\textit{#1}}
\providecommand{\KeywordTok}[1]{\textbf{#1}}
\providecommand{\NormalTok}[1]{#1}
\providecommand{\StringTok}[1]{\textit{#1}}
\providecommand{\OperatorTok}[1]{#1}
\providecommand{\DataTypeTok}[1]{#1}
\providecommand{\FunctionTok}[1]{#1}

% --- Links & references ---
\usepackage[hidelinks,breaklinks=true]{hyperref}
\usepackage{cleveref}

% --- Bibliography (APA7 via natbib + apalike) ---
% APJM/MIR accept APA7. natbib's apalike close approximation; for full APA7
% compliance use biblatex with apa style.
\usepackage[round,sort]{natbib}
\bibliographystyle{apalike}

% --- Pandoc compatibility shims ---
\providecommand{\tightlist}{\setlength{\itemsep}{0pt}\setlength{\parskip}{0pt}}
\providecommand{\passthrough}[1]{\lstset{mathescape=false}#1\lstset{mathescape=true}}

% --- Title page (filled in per paper) ---
\title{p3\_vietnam (bản dịch tiếng Việt)}
\author{Author details withheld for blind review}
\date{2026-05-21}

\begin{document}

\maketitle

\hypertarget{xem-xuxe9t-lux1ea1i-mux1ed1i-quan-hux1ec7-giux1eefa-quux1ed1c-tux1ebf-houxe1-vuxe0-hiux1ec7u-quux1ea3-houx1ea1t-ux111ux1ed9ng-kinh-doanh-cux1ee7a-doanh-nghiux1ec7p-tux1ea1i-mux1ed9t-thux1ecb-trux1b0ux1eddng-mux1edbi-nux1ed5i-vai-truxf2-cux1ee7a-nux103ng-lux1ef1c-cuxf4ng-nghux1ec7-vuxe0-mux1ee9c-ux111ux1ed9-uxe1p-dux1ee5ng-sux1ed1-nux1ec1n-tux1ea3ng}{%
\section{Xem xét lại mối quan hệ giữa quốc tế hoá và hiệu quả hoạt động
kinh doanh của doanh nghiệp tại một thị trường mới nổi: Vai trò của năng
lực công nghệ và mức độ áp dụng số nền
tảng}\label{xem-xuxe9t-lux1ea1i-mux1ed1i-quan-hux1ec7-giux1eefa-quux1ed1c-tux1ebf-houxe1-vuxe0-hiux1ec7u-quux1ea3-houx1ea1t-ux111ux1ed9ng-kinh-doanh-cux1ee7a-doanh-nghiux1ec7p-tux1ea1i-mux1ed9t-thux1ecb-trux1b0ux1eddng-mux1edbi-nux1ed5i-vai-truxf2-cux1ee7a-nux103ng-lux1ef1c-cuxf4ng-nghux1ec7-vuxe0-mux1ee9c-ux111ux1ed9-uxe1p-dux1ee5ng-sux1ed1-nux1ec1n-tux1ea3ng}}

\emph{Thông tin tác giả đã được lược bỏ để phục vụ bình duyệt ẩn danh}

\emph{Ngày 6 tháng 5 năm 2026}

\emph{Tạp chí mục tiêu: Asia Pacific Journal of Management (APJM)}

\textbf{Phân loại bản thảo:} bài báo nghiên cứu.

\textbf{Số từ} (phần thân chính, không tính tóm tắt, tài liệu tham khảo,
bảng và hình): khoảng 6.800 từ.

\textbf{Bảng:} 4 (Bảng 1 thống kê mô tả; Bảng 2 tổng hợp hệ số trọng
tâm; Bảng 3 các bảng kiểm tra độ vững; Bảng 4 các điểm ngoặt).

\textbf{Hình:} 6 (Hình 1 mô hình khái niệm; các Hình 2a--2d đường cong
FSTS theo từng đợt khảo sát; Hình 3 đường cong điều tiết tại p25 so với
p75 của biến điều tiết).

\begin{center}\rule{0.5\linewidth}{0.5pt}\end{center}

\hypertarget{tuxf3m-tux1eaft}{%
\subsection{Tóm tắt}\label{tuxf3m-tux1eaft}}

\textbf{Mục đích.} Nghiên cứu xem xét lại mối quan hệ giữa quốc tế hoá
và hiệu quả hoạt động kinh doanh của doanh nghiệp tại một thị trường mới
nổi, đồng thời khảo sát vai trò của năng lực công nghệ nước ngoài và áp
dụng công nghệ số nền tảng đối với năng suất doanh nghiệp trong điều
kiện chuyển đổi thể chế và chuyển đổi số. Lấy Việt Nam làm trọng tâm,
nghiên cứu đặt câu hỏi liệu cường độ xuất khẩu có liên hệ phi tuyến với
năng suất lao động hay không, và liệu hai cấu trúc năng lực này có đóng
những vai trò khác biệt theo từng đợt khảo sát hay không.

\textbf{Thiết kế / phương pháp luận / cách tiếp cận.} Nghiên cứu sử dụng
ba đợt dữ liệu vi mô từ Khảo sát Doanh nghiệp của Ngân hàng Thế giới cho
Việt Nam (2009, 2015 và 2023; N gộp = 2.958), ước lượng các mô hình bình
phương nhỏ nhất theo từng đợt và gộp, kèm sai số chuẩn vững HC1, số hạng
bậc hai của cường độ xuất khẩu và các đặc tả tương tác. Chỉ số Năng lực
Công nghệ (Technological Capability Index, TCI) được xây dựng từ chứng
nhận chất lượng được công nhận quốc tế và công nghệ được cấp phép nước
ngoài; Chỉ số Áp dụng Số nền tảng (Digital Adoption Index, DAI) đo bằng
sự hiện diện website. Hai chỉ số không chia sẻ bất kỳ chỉ báo nào.

\textbf{Kết quả nghiên cứu.} Lợi ích năng suất của quốc tế hoá tập trung
tại biên tham gia, tức bước chuyển từ không xuất khẩu sang xuất khẩu,
trong khi cường độ xuất khẩu trong nhóm doanh nghiệp đã xuất khẩu mang
lại lợi suất biên giảm dần đến triệt tiêu. Đường chữ U ngược toàn mẫu
(điểm ngoặt 39-46\% cường độ xuất khẩu) được nhận diện chủ yếu thông qua
bước tham gia này chứ không qua độ cong dày đặc bên trong nhóm doanh
nghiệp xuất khẩu. Năng lực công nghệ thể hiện liên hệ thuận chiều với
năng suất trong cả ba đợt; áp dụng số nền tảng có hệ số dương trong 2009
và 2023 nhưng không ý nghĩa thống kê trong 2015, và ước lượng biến công
cụ cho thấy liên hệ thuận chiều của DAI ước lượng bằng OLS nhiều khả
năng phản ánh chọn lựa hơn là quan hệ nhân quả.

\textbf{Tính mới / đóng góp.} Nghiên cứu đóng góp cho tài liệu về thị
trường mới nổi trong kinh doanh quốc tế bằng cách phân tách năng lực
công nghệ nước ngoài khỏi áp dụng số nền tảng và bằng cách cho thấy các
tác động số gộp có thể che lấp tính không đồng nhất theo thời gian đáng
kể. Phát hiện gợi ý rằng tính liên quan đến năng suất của áp dụng số cơ
bản nhạy cảm với bối cảnh và đặc trưng theo đợt khảo sát hơn là ổn định
đồng nhất qua các giai đoạn quốc tế hoá.

\textbf{Từ khoá:} quốc tế hoá-hiệu quả hoạt động; thị trường mới nổi; áp
dụng công nghệ số; năng lực công nghệ; Việt Nam; năng suất doanh nghiệp.

\textbf{Phân loại JEL:} F23; O33; D22; L25; O53.

\textbf{Loại bài:} Bài báo nghiên cứu.

\hypertarget{ux111iux1ec3m-nhux1ea5n}{%
\subsection{Điểm nhấn}\label{ux111iux1ec3m-nhux1ea5n}}

\begin{itemize}
\item
  Mối quan hệ giữa quốc tế hoá và hiệu quả hoạt động ở doanh nghiệp Việt
  Nam có tính phi tuyến vững chắc: kiểm định Lind--Mehlum (U-shape test)
  bác bỏ tính đơn điệu trong cả ba đợt (2009 p = .006, 2015 p = .009,
  2023 p = .013) và trong mẫu gộp (p \textless{} .001), với các điểm
  ngoặt tập trung trong khoảng 39 đến 46 \% cường độ xuất khẩu trực
  tiếp.
\item
  Năng lực công nghệ và áp dụng công nghệ số được tách biệt như những
  thước đo chính không chồng lấn: \(\tilde{\mathrm{TCI}}\) là trung bình
  chuẩn hoá theo z (z-standardisation) trong từng đợt của chứng nhận
  chất lượng và công nghệ được cấp phép nước ngoài, trong khi
  \(\tilde{\mathrm{DAI}}\) là chỉ báo hiện diện số nền tảng dựa trên
  website đã chuẩn hoá theo z trong từng đợt.
\item
  \(\tilde{\mathrm{TCI}}\) có liên hệ thuận chiều với năng suất trong cả
  ba đợt (β = 0,215; 0,128; 0,123) và trong mẫu gộp (β = 0,179, p
  \textless{} .001), và điều tiết độ cong trong ba trong bốn bảng (M3
  joint p = .040, .713, .027 và .003).
\item
  \(\tilde{\mathrm{DAI}}\) mạnh nhất vào năm 2009 (β = 0,175, p
  \textless{} .001), không có ý nghĩa năm 2015 (β = -0,044, p = .377),
  dương năm 2023 (β = 0,095, p = .038) và trong mẫu gộp (β = 0,078, p =
  .004); các dịch chuyển giữa các đợt có thể phân biệt được về mặt thống
  kê (kiểm định z Paternoster z = 3,353 giữa 2009 và 2015; z = -2,051
  giữa 2015 và 2023).
\item
  Điều tiết của DAI tập trung vào năm 2023, trong đó
  \(\mathrm{FSTS}^{c}\) × \(\tilde{\mathrm{DAI}}\) = -0,912 (p = .043)
  và kiểm định liên hợp M8 ở mức cận ngưỡng (p = .062); kiểm định liên
  hợp M8 gộp cũng ở mức cận ngưỡng (p = .083), được chi phối bởi năm
  2023 chứ không phải bởi một điều tiết ổn định xuyên suốt giai đoạn.
\end{itemize}

\hypertarget{1-giux1edbi-thiux1ec7u}{%
\subsection{1. Giới thiệu}\label{1-giux1edbi-thiux1ec7u}}

\hypertarget{11-bux1ed1i-cux1ea3nh-vuxe0-ux111ux1ed9ng-cux1a1-nghiuxean-cux1ee9u}{%
\subsubsection{1.1 Bối cảnh và động cơ nghiên
cứu}\label{11-bux1ed1i-cux1ea3nh-vuxe0-ux111ux1ed9ng-cux1a1-nghiuxean-cux1ee9u}}

Việt Nam là bối cảnh có giá trị phân tích để xem xét lại mối quan hệ
giữa quốc tế hoá và hiệu quả hoạt động kinh doanh của doanh nghiệp
(I--P). Tại đây, các doanh nghiệp mở rộng ra nước ngoài trong điều kiện
chuyển đổi thể chế, tích luỹ năng lực không đồng đều và hạ tầng số biến
đổi nhanh chóng.

Trong bối cảnh như vậy, không nên giả định rằng quốc tế hoá sẽ tạo ra
khoản phụ trội hiệu quả tuyến tính đơn giản. Doanh nghiệp có thể tiếp
cận các thị trường lớn hơn, hưởng lợi từ học hỏi và đa dạng hoá dòng
doanh thu. Tuy nhiên, doanh nghiệp cũng có thể đối mặt với chi phí phối
hợp (coordination cost) tăng cao, gánh nặng xử lý thông tin và sức ép tổ
chức khi mức độ tham gia nước ngoài sâu hơn (Wright et al., 2005;
Cuervo-Cazurra and Genc, 2008; Wu et al., 2016).

Sự căng thẳng này nằm ở trung tâm của tài liệu I--P. Một truyền thống
nghiên cứu lâu đời lập luận rằng quốc tế hoá có thể cải thiện hiệu quả ở
các mức thấp và trung gian thông qua quy mô, học hỏi và đa dạng hoá,
đồng thời tạo ra lợi suất giảm dần hoặc âm ở các mức cao hơn do tính
phức tạp và gánh nặng phối hợp. Bằng chứng từ phân tích tổng hợp mạnh mẽ
ủng hộ quan điểm rằng tính phi tuyến là một đặc trưng trung tâm của mối
quan hệ này chứ không phải một bất thường thực nghiệm (Vernon, 1979; Lu
and Beamish, 2004; Hennart, 2007; Coviello et al., 2017; Marano et al.,
2016). Bằng chứng cấp doanh nghiệp gần đây từ Việt Nam xác nhận vòng
xoáy năng suất-xuất khẩu hai chiều ở các SME (Nguyen và cộng sự, 2025;
Dang và cộng sự, 2024) và nhấn mạnh chuyển đổi số như một động lực mạnh
cho hiệu quả xuất khẩu (Tran và cộng sự, 2025; Ngo và cộng sự, 2024).

Số hoá bổ sung thêm một lớp phức tạp cho cuộc tranh luận này. Công cụ số
có thể giảm ma sát giao tiếp, đẩy nhanh giao dịch và hỗ trợ phối hợp
xuyên biên giới. Tuy nhiên, những lợi ích đó không phát sinh tự động.
Giá trị thực tế của chúng phụ thuộc vào ba điều kiện ở doanh nghiệp:
chiều sâu tổ chức, năng lực hấp thụ (absorptive capacity), và các thông
lệ bổ trợ cần thiết để chuyển hoá việc áp dụng công nghệ số thành lợi
ích năng suất (Cohen and Levinthal, 1990; Vial, 2019; Verhoef et al.,
2021; Stallkamp and Schotter, 2021; Petricevic and Teece, 2019). Vì lý
do này, không nên xem năng lực số như nguồn lực có lợi phổ quát với mức
sinh lời không đổi giữa các doanh nghiệp và theo thời gian.

Bối cảnh Việt Nam khiến vấn đề này trở nên đặc biệt quan trọng. Vì doanh
nghiệp được nhúng trong một nền kinh tế đang nâng cấp năng lực và chuyển
đổi số, lợi ích từ quốc tế hoá và mức sinh lời từ số hoá đều có thể thay
đổi giữa các giai đoạn phát triển. Điều này đặt ra một câu hỏi trung tâm
cho nghiên cứu hiện tại: liệu năng lực số ở Việt Nam có vận hành như một
tài sản ổn định nâng cao hiệu quả, hay nó vận hành như một nguồn lực phụ
thuộc giai đoạn (stage-contingent) với giá trị thay đổi theo vòng đời
quốc tế hoá?

Ba điểm ngoặt thể chế định hình cửa sổ quan sát. Việt Nam gia nhập Tổ
chức Thương mại Thế giới đầu năm 2007, mở ra giai đoạn trước đợt khảo
sát 2009 và chuyển một nhóm doanh nghiệp xuất khẩu định hướng nội địa
thành nhóm có độ tiếp xúc thị trường rộng hơn nhưng hạ tầng hấp thụ còn
hạn chế. Đợt 2015 ghi lại giai đoạn giữa của pha thứ hai, khi xuất khẩu
sản xuất mở rộng cùng tồn tại với hạ tầng thương mại số kém phát triển,
tích hợp logistics xuyên biên giới yếu, và một nhóm doanh nghiệp xuất
khẩu vẫn tập trung ở các phân khúc thâm dụng lao động. Đợt 2023 diễn ra
sau ba diễn biến chính: Chương trình Chuyển đổi Số Quốc gia khởi động
năm 2020; sự mở rộng nhanh chóng của các nền tảng thanh toán điện tử và
thương mại điện tử xuyên biên giới; và tái cân bằng đầu tư trực tiếp
nước ngoài (FDI) hướng tới sản xuất qua trung gian số và liên kết dịch
vụ. Ba đợt khảo sát do đó quan sát doanh nghiệp dưới những tổ hợp khác
biệt về cấu trúc giữa áp lực quốc tế hoá và mức độ sẵn có của hạ tầng
số. Đây chính là điều khiến cách đọc theo vòng đời có thể kiểm định
được, chứ không thuần tuý khái niệm.

Những dịch chuyển này không mang tính hình thức. Bản thân thành phần của
nhóm doanh nghiệp xuất khẩu cũng tiến triển qua ba đợt. Tỉ lệ doanh
nghiệp báo cáo cường độ xuất khẩu trực tiếp dương suy giảm khi các doanh
nghiệp dịch vụ và doanh nghiệp chuỗi cung ứng liên kết FDI gia nhập mẫu.
Cùng lúc, mức trung bình của áp dụng công nghệ số nền tảng tăng lên cùng
sự lan toả của website, hệ thống thanh toán điện tử và các giao diện
giao dịch số. Đọc mối quan hệ I--P và kênh áp dụng công nghệ số như
những sự kiện cấu trúc cố định xuyên suốt cửa sổ 14 năm này sẽ bỏ qua
một thực tế: quần thể doanh nghiệp nền, các chi phí phối hợp ràng buộc
và bộ khung thể chế cho thương mại xuyên biên giới đều thay đổi đáng kể
về mặt vật chất.

\hypertarget{12-khoux1ea3ng-trux1ed1ng-nghiuxean-cux1ee9u}{%
\subsubsection{1.2 Khoảng trống nghiên
cứu}\label{12-khoux1ea3ng-trux1ed1ng-nghiuxean-cux1ee9u}}

Ba khoảng trống thúc đẩy nghiên cứu này.

Thứ nhất, các nghiên cứu hiện hữu thường xem số hoá như một nguồn lực
tích cực một cách rộng rãi mà chưa chú ý đầy đủ đến sự biến đổi theo
thời gian và bối cảnh trong mức sinh lời của nó. Cách xử lý như vậy có
nguy cơ ngụ ý rằng áp dụng công nghệ số tạo ra một khoản phụ trội hiệu
quả tương đối ổn định giữa các doanh nghiệp và qua các giai đoạn chuyển
đổi kinh tế (Strange and Zucchella, 2017; Goldfarb and Tucker, 2019).
Trên thực tế, số hoá có thể sinh ra lợi suất không đồng đều bởi vì các
doanh nghiệp khác biệt về quy mô, thông lệ và năng lực bổ trợ.

Thứ hai, sự phân biệt giữa năng lực công nghệ và áp dụng công nghệ số
vẫn chưa được phát triển đầy đủ. Năng lực công nghệ đề cập đến những trữ
lượng học hỏi, giải quyết vấn đề, cải tiến quy trình và năng lực đổi mới
sâu hơn bên trong doanh nghiệp (Lall, 1992). Ngược lại, áp dụng công
nghệ số nền tảng phản ánh một lớp cơ bản hơn về mức độ sẵn sàng số và
các giao diện hoặc cơ chế giao dịch có hỗ trợ số (Bharadwaj et al.,
2013; Verhoef et al., 2021; Hanelt et al., 2021; Brouthers et al.,
2016). Mặc dù hai khái niệm này có liên quan, không nên xem chúng có thể
hoán đổi cho nhau. Việc gộp chúng vào một biến tổng có thể làm giảm tính
rõ ràng của khái niệm và làm mờ các cơ chế gắn số hoá với hiệu quả.

Thứ ba, ước lượng gộp có thể che giấu sự không đồng nhất theo vòng đời
đáng kể. Trong bối cảnh như Việt Nam, nơi doanh nghiệp và thể chế đi qua
các giai đoạn chuyển đổi khác biệt, vai trò của quốc tế hoá, năng lực
công nghệ và áp dụng công nghệ số có thể biến đổi đáng kể giữa các đợt
khảo sát. Mô hình gộp hữu ích để xác định các tác động trung bình, nhưng
không đủ nếu các mối quan hệ nền dịch chuyển theo thời gian. Do đó, một
thiết kế kết hợp phân tích gộp và phân tích theo từng đợt là cần thiết
để xác định liệu các tác động quan sát được ổn định hay phụ thuộc giai
đoạn.

Nghiên cứu phân tích tổng hợp gần đây củng cố mối quan tâm này. Pisani,
Garcia-Bernardo, và Heemskerk (2020, SMJ) chỉ ra rằng mối quan hệ chữ U
ngược giữa quốc tế hoá và hiệu quả hoạt động yếu đi đáng kể (và đôi khi
biến mất) dưới các thiết kế nhận diện chặt chẽ hơn trên các mẫu gộp đa
quốc gia. Phát hiện này gợi ý rằng việc tổng hợp giữa các bối cảnh thể
chế khác biệt làm gia tăng tính đều đặn về mặt dạng hàm bề ngoài. Wu,
Fan, và Chen (2022, MIR) mở rộng lập luận này. Trong phân tích tổng hợp
về doanh nghiệp đa quốc gia thị trường mới nổi (EMNE), họ cho thấy bối
cảnh thể chế điều tiết quy mô tác động I--P mạnh mẽ hơn các biến năng
lực ở cấp doanh nghiệp. Tóm lại, những phát hiện này ngụ ý rằng các
thiết kế dọc trong nội bộ quốc gia (giữ nguyên bối cảnh thể chế trong
khi để thời gian biến đổi) cung cấp kiểm định đáng tin cậy hơn. Cụ thể,
kiểm định này đánh giá liệu chữ U ngược là hiện tượng giả do gộp đa quốc
gia hay là mối quan hệ thực sự tồn tại khi môi trường thể chế tiến hoá.
Bảng dữ liệu WBES ba đợt của Việt Nam (2009, 2015, 2023) cung cấp đúng
thiết kế này.

\hypertarget{13-ux111uxf3ng-guxf3p}{%
\subsubsection{1.3 Đóng góp}\label{13-ux111uxf3ng-guxf3p}}

Nghiên cứu này có ba đóng góp vào tài liệu.

Thứ nhất, nó tinh chỉnh cuộc tranh luận I--P bằng cách cho thấy bằng
chứng từ Việt Nam ủng hộ mối quan hệ phi tuyến, nhưng tính nổi bật và
khả năng quan sát của mối quan hệ đó biến đổi theo thời gian. Thay vì
xem tính phi tuyến như một sự kiện cấu trúc cố định xuất hiện giống nhau
trong mọi giai đoạn, phân tích cho thấy đường cong I--P phải được đọc
cùng với môi trường năng lực rộng hơn.

Thứ hai, nghiên cứu cải thiện tính giá trị khái niệm bằng cách tách năng
lực công nghệ khỏi áp dụng công nghệ số nền tảng. Sự phân biệt này quan
trọng về mặt lý thuyết vì các trữ lượng năng lực sâu hơn và mức kích
hoạt số cơ bản có thể tạo ra hiệu quả qua các kênh khác nhau. Nó cũng
quan trọng về mặt thực nghiệm vì hai khái niệm không thể hiện các mẫu
hình giống nhau qua các đợt khảo sát ở Việt Nam.

Thứ ba, nghiên cứu giới thiệu một cách diễn giải vòng đời về quốc tế hoá
số. Bằng chứng gợi ý rằng năng lực số không phải là một khoản phụ trội
ổn định phổ quát cũng không phải là một nguồn lực không hiệu quả đồng
nhất. Thay vào đó, đó là một nguồn không đồng đều và phụ thuộc giai đoạn
của tính không đồng nhất hiệu quả. Quan điểm này giúp giải thích vì sao
các tác động trung bình gộp có thể cùng tồn tại với những khác biệt đáng
kể theo từng đợt.

Đóng góp thứ tư mang tính bối cảnh: mẫu hình thực nghiệm của Việt Nam
cung cấp một phân rã tham gia-cường độ. Cụ thể, chữ U ngược được nhận
diện vững chắc (điểm ngoặt 39--46 \% FSTS) chủ yếu được dẫn dắt bởi biên
tham gia. Trong nhóm doanh nghiệp xuất khẩu, biến thiên cường độ mang
lại lợi suất phẳng đến giảm dần (Panel H), đặt giới hạn cho cách đọc
thông thường về mẫu hình chữ U ngược. Sự khác biệt giữa mẫu hình này và
những gì được ghi nhận ở các nền kinh tế tiên tiến về số (nơi cơ chế chi
phí phối hợp được làm dịu bởi hạ tầng số toàn diện) hướng tới các cơ chế
cấp thể chế như những biến điều tiết ứng viên về vị trí đường cong I--P.
Bằng chứng nội bối cảnh này neo vào phần chuyển tiếp số của phổ thể chế.

\hypertarget{14-bux1ed1-cux1ee5c}{%
\subsubsection{1.4 Bố cục}\label{14-bux1ed1-cux1ee5c}}

Phần còn lại của bài báo được tổ chức như sau. Mục 2 phát triển khung lý
thuyết và các giả thuyết. Mục 3 mô tả dữ liệu, biến số và chiến lược
thực nghiệm. Mục 4 trình bày kết quả. Mục 5 thảo luận hàm ý lý thuyết và
quản trị. Mục 6 kết luận.

\hypertarget{2-cux1a1-sux1edf-luxfd-thuyux1ebft-vuxe0-giux1ea3-thuyux1ebft}{%
\subsection{2. Cơ sở lý thuyết và giả
thuyết}\label{2-cux1a1-sux1edf-luxfd-thuyux1ebft-vuxe0-giux1ea3-thuyux1ebft}}

\hypertarget{21-quux1ed1c-tux1ebf-houxe1-vuxe0-hiux1ec7u-quux1ea3-houx1ea1t-ux111ux1ed9ng-kinh-doanh-cux1ee7a-doanh-nghiux1ec7p}{%
\subsubsection{2.1 Quốc tế hoá và hiệu quả hoạt động kinh doanh của
doanh
nghiệp}\label{21-quux1ed1c-tux1ebf-houxe1-vuxe0-hiux1ec7u-quux1ea3-houx1ea1t-ux111ux1ed9ng-kinh-doanh-cux1ee7a-doanh-nghiux1ec7p}}

Mối quan hệ giữa quốc tế hoá và hiệu quả hoạt động kinh doanh của doanh
nghiệp khó có thể là tuyến tính trong một nền kinh tế chuyển đổi như
Việt Nam. Mở rộng ra nước ngoài có thể cải thiện hiệu quả thông qua phạm
vi thị trường lớn hơn, đa dạng hoá và học hỏi từ các môi trường bên
ngoài. Doanh nghiệp có thể sử dụng quốc tế hoá để phân bổ chi phí cố
định, tiếp cận khách hàng mới và thu nhận kiến thức hỗ trợ cải tiến vận
hành.

Đồng thời, mức độ tham gia quốc tế sâu hơn thường tạo ra chi phí phối
hợp và gánh nặng tổ chức. Doanh nghiệp phải quản lý nhiều thị trường,
dung hoà các yêu cầu khách hàng đa dạng, xử lý thông tin nhiều hơn và
duy trì kiểm soát quản trị lớn hơn. Khi mở rộng nước ngoài tăng cường,
các chi phí này có thể tăng nhanh hơn lợi ích, tạo ra lợi suất giảm dần
hoặc thậm chí suy giảm hiệu quả. Logic cơ bản này làm nền tảng cho quan
điểm phi tuyến cổ điển về mối quan hệ I--P (Contractor, 2007; Hennart,
2011; Marano et al., 2016). Các cách giải thích theo quá trình về quốc
tế hoá, bao gồm khung Uppsala đã cập nhật, cũng nhấn mạnh rằng lợi ích
về hiệu quả và chi phí điều chỉnh diễn ra dần dần khi doanh nghiệp tích
luỹ kiến thức và cam kết thị trường (Vahlne and Johanson, 2017; Knight
and Liesch, 2016).

Lập luận này đặc biệt khả tín ở Việt Nam vì doanh nghiệp hoạt động dưới
điều kiện năng lực không đồng đều. Một số doanh nghiệp có thể chuyển hoá
mở rộng nước ngoài thành lợi thế học hỏi và quy mô, trong khi những
doanh nghiệp khác có thể gặp sức ép tổ chức sớm hơn. Kỳ vọng phù hợp do
đó không phải là một độ dốc dương đồng nhất, mà là một mối quan hệ phi
tuyến trong đó lợi ích từ quốc tế hoá trở nên ngày càng khó duy trì.

Hai lực đối lập làm nền cho độ cong này.

Ở mặt tích cực, cường độ xuất khẩu trực tiếp tăng tạo ra lợi thế kinh tế
theo quy mô, lan toả tri thức từ khách hàng nước ngoài, và tác động
học-bằng-xuất-khẩu giúp nâng năng suất (Wagner, 2007).

Ở mặt tiêu cực, phối hợp sản xuất cho các thị trường có khoảng cách thể
chế áp đặt các chi phí xử lý thông tin tăng phi tuyến. Mỗi thị trường
nước ngoài bổ sung thêm các yêu cầu tuân thủ, gánh nặng quan hệ khách
hàng và sự phụ thuộc chuỗi cung ứng. Chi phí phối hợp biên tăng nhanh
hơn lợi ích quy mô biên khi vượt qua ngưỡng nhất định. Trong bối cảnh
chuyển đổi nơi cảng, định chế tài trợ thương mại (trade finance), thị
trường số và cơ chế giải quyết tranh chấp vẫn đang trưởng thành, ngưỡng
này có thể ràng buộc ở mức cường độ xuất khẩu thấp hơn so với các nền
kinh tế trưởng thành. Đặc điểm này làm sắc nét hơn độ cong so với đường
cơ sở phân tích tổng hợp (Hennart, 2007; Wagner, 2007; Marano et al.,
2016). Cơ chế là chi phí giao dịch thể chế (Williamson, 1985; Hennart,
2007): trong nền kinh tế chuyển đổi, các lỗ hổng thực thi, thâm hụt
logistics và bất đối xứng thông tin làm tăng chi phí biên của việc quản
lý mỗi thị trường nước ngoài bổ sung. Vì vậy, sự suy giảm sau ngưỡng
ràng buộc ở mức cường độ xuất khẩu thấp hơn so với các bối cảnh nơi chất
lượng hạ tầng số và thực thi hợp đồng đã hấp thụ phần lớn ma sát này.

Một cân nhắc thứ hai phát sinh từ cấu trúc của cường độ xuất khẩu trực
tiếp trong WBES Việt Nam: biến FSTS bị chặn ở mức không và bị lạm phát ở
mức không nghiêm trọng (tỉ lệ doanh nghiệp không xuất khẩu là 71,6 \%
năm 2009, 79,3 \% năm 2015 và 81,2 \% năm 2023; gộp 77,4 \%). Mối quan
hệ quốc tế hoá--hiệu quả do đó bao gồm hai biên có thể phân biệt về mặt
phân tích. Biên tham gia là bước từ FSTS = 0 sang FSTS \textgreater{} 0,
nơi doanh nghiệp vượt qua ngưỡng nhập cuộc thị trường quốc tế và hấp thụ
chi phí nhập cuộc cố định. Biên cường độ là biến thiên của FSTS trong
nhóm doanh nghiệp xuất khẩu.

Lý thuyết dự đoán rằng lợi ích năng suất có thể tích luỹ ở cả hai biên:
biên tham gia ghi nhận tác động học-bằng-xuất-khẩu và sự tiếp xúc với
nhu cầu thị trường nước ngoài; biên cường độ ghi nhận lợi thế kinh tế
theo quy mô ở mặt tích cực và gánh nặng phối hợp ở mặt tiêu cực (Bernard
et al., 2007; Wagner, 2007). Trong các nền kinh tế chuyển đổi với hạ
tầng xuất khẩu mỏng, biên tham gia khả năng là biên nâng cao năng suất
chủ đạo, trong khi biên cường độ có thể cho thấy lợi suất giảm dần hoặc
triệt tiêu khi đã vượt qua biên tham gia.

\textbf{Giả thuyết H1.} Mối quan hệ quốc tế hoá--hiệu quả ở Việt Nam là
không đơn điệu trong toàn mẫu và được hiểu tốt nhất qua một cấu trúc
tham gia-cường độ trong bối cảnh xuất khẩu lạm phát ở mức không.
\textbf{(H1a, biên tham gia)} Việc chuyển từ không xuất khẩu (FSTS = 0)
sang xuất khẩu (FSTS \textgreater{} 0) có liên hệ thuận chiều với năng
suất lao động, phản ánh bước nhảy về học hỏi và quy mô khi nhập cuộc.
\textbf{(H1b, biên cường độ)} Trong nhóm doanh nghiệp xuất khẩu, cường
độ xuất khẩu trực tiếp bổ sung được kỳ vọng cho thấy lợi suất biên yếu
hơn, giảm dần hoặc không có ý nghĩa so với biên tham gia.

Đặc tả chỉ riêng doanh nghiệp xuất khẩu (Panel H) cung cấp kiểm định
trực tiếp nhất cho H1b: nếu độ cong chữ U ngược phản ánh biến thiên
cường độ trong nhóm doanh nghiệp xuất khẩu, số hạng bậc hai phải còn ý
nghĩa trong nhóm phụ này. Mẫu hình gần phẳng quan sát được trong Panel H
(gộp \((\mathrm{FSTS}^{c})^2\) β = −0,200, p = .660) thay vào đó cho
thấy biên tham gia (H1a) là nguồn chủ đạo của độ cong toàn mẫu.

\hypertarget{22-nux103ng-lux1ef1c-cuxf4ng-nghux1ec7-nux1b0ux1edbc-ngouxe0i-vuxe0-tiuxeau-chuux1ea9n-cuxf9ng-hiux1ec7u-quux1ea3-houx1ea1t-ux111ux1ed9ng-kinh-doanh-cux1ee7a-doanh-nghiux1ec7p}{%
\subsubsection{2.2 Năng lực công nghệ nước ngoài và tiêu chuẩn cùng hiệu
quả hoạt động kinh doanh của doanh
nghiệp}\label{22-nux103ng-lux1ef1c-cuxf4ng-nghux1ec7-nux1b0ux1edbc-ngouxe0i-vuxe0-tiuxeau-chuux1ea9n-cuxf9ng-hiux1ec7u-quux1ea3-houx1ea1t-ux111ux1ed9ng-kinh-doanh-cux1ee7a-doanh-nghiux1ec7p}}

Theo truyền thống Lall (1992) cho các doanh nghiệp thị trường mới nổi,
bài báo này sử dụng cách đọc chặt chẽ về đo lường đối với năng lực công
nghệ. Cụ thể, năng lực công nghệ nước ngoài và tiêu chuẩn ghi nhận sự
tiếp xúc của doanh nghiệp với các đầu vào công nghệ được xác nhận từ bên
ngoài: chứng nhận chất lượng được công nhận quốc tế và công nghệ được
cấp phép nước ngoài. Đây là một mặt có thể quan sát của khái niệm năng
lực hấp thụ rộng hơn của Cohen và Levinthal (1990) và khái niệm năng lực
động của Teece (2007). Tuy nhiên, thước đo hiện tại không tuyên bố nhận
diện toàn bộ trữ lượng năng lực hấp thụ; nó chỉ ghi nhận thành phần
hướng ngoại của trữ lượng đó (Eisenhardt and Martin, 2000; Karna et al.,
2016).

Trong bối cảnh quốc tế, năng lực hướng ngoại này đặc biệt quan trọng bởi
vì doanh nghiệp phải ứng phó với các thị trường xa lạ, đáp ứng tiêu
chuẩn chất lượng nước ngoài và tích hợp công nghệ được cấp phép từ bên
ngoài vào các thông lệ tổ chức.

Một doanh nghiệp với năng lực công nghệ nước ngoài / tiêu chuẩn mạnh hơn
có nhiều khả năng chuyển hoá sự tiếp xúc quốc tế thành lợi ích năng
suất. Doanh nghiệp có thể điều chỉnh sản phẩm và quy trình để đáp ứng
các yêu cầu nước ngoài hiệu quả hơn, tích hợp công nghệ nước ngoài được
cấp phép vào sản xuất, và đối phó tốt hơn với các yêu cầu vận hành phát
sinh từ hoạt động xuất khẩu. Ngay cả khi năng lực này không thay đổi cơ
bản độ cong của mối quan hệ I--P, nó vẫn nên cải thiện mức độ hiệu quả
tổng thể của doanh nghiệp bằng cách nâng sàn năng suất giữa các doanh
nghiệp đã tiếp xúc.

Điều này ngụ ý một mối liên hệ thuận chiều trực tiếp giữa năng lực công
nghệ nước ngoài / tiêu chuẩn và hiệu quả hoạt động kinh doanh của doanh
nghiệp ở Việt Nam. Chúng tôi xem nó ở đây như một khái niệm nên nâng
năng lực hưởng lợi từ quốc tế hoá của doanh nghiệp và cũng nên hỗ trợ
năng suất rộng hơn, mà không giả định rằng nó đại diện cho chiều đầy đủ
về đổi mới và R\&D.

Về mặt vận hành, \(\tilde{\mathrm{TCI}}\) chính được xây dựng từ hai chỉ
báo: chứng nhận chất lượng được công nhận quốc tế (b8) và công nghệ được
cấp phép nước ngoài (e6). Các chỉ báo này đo lường mức độ tiếp xúc với
các kênh công nghệ và tiêu chuẩn nước ngoài chứ không phải nỗ lực R\&D
nội bộ hay hoạt động bằng sáng chế. Một chỉ số hợp thành cấu thành
(formative composite) mở rộng đổi mới (\(\mathrm{TCI}_{\text{full}}\),
bổ sung đổi mới sản phẩm h1 và hoạt động R\&D h8) được báo cáo trong mục
4.5 Panel A như điều kiện biên. Nếu các tác động trực tiếp suy giảm khi
các chỉ báo đổi mới được bổ sung, điều đó cho thấy thước đo chính cung
cấp thông tin cụ thể về kênh công nghệ và tiêu chuẩn nước ngoài, chứ
không phải về trữ lượng năng lực hấp thụ rộng hơn.

\textbf{Giả thuyết H2.} Năng lực công nghệ nước ngoài / tiêu chuẩn
(\(\tilde{\mathrm{TCI}}\)) có liên hệ thuận chiều với hiệu quả hoạt động
kinh doanh của doanh nghiệp ở Việt Nam.

\hypertarget{23-hiux1ec7n-diux1ec7n-sux1ed1-dux1ef1a-truxean-website-vuxe0-hiux1ec7u-quux1ea3-houx1ea1t-ux111ux1ed9ng-kinh-doanh-cux1ee7a-doanh-nghiux1ec7p}{%
\subsubsection{2.3 Hiện diện số dựa trên website và hiệu quả hoạt động
kinh doanh của doanh
nghiệp}\label{23-hiux1ec7n-diux1ec7n-sux1ed1-dux1ef1a-truxean-website-vuxe0-hiux1ec7u-quux1ea3-houx1ea1t-ux111ux1ed9ng-kinh-doanh-cux1ee7a-doanh-nghiux1ec7p}}

\(\tilde{\mathrm{DAI}}\) chính được sử dụng trong bài báo này là một
thước đo hiện diện số dựa trên website: một chỉ báo nhị phân về việc
doanh nghiệp có sở hữu website riêng hay không. Đây là một dấu mốc nền
tảng và có thể so sánh giữa các đợt về áp dụng công nghệ số, nó không đo
lường tích hợp số ở cấp độ giao dịch, hạ tầng thanh toán điện tử, hay
chuyển đổi số theo nghĩa của Bharadwaj et al. (2013) / Verhoef et al.
(2021) / Vial (2019). Chúng tôi sử dụng nó ở đây chính vì đó là chỉ báo
số duy nhất mà công cụ WBES mang theo có thể so sánh được giữa các đợt
2009 / 2015 / 2023; các chỉ báo cấp giao dịch như tỉ trọng thanh toán
điện tử chỉ xuất hiện trong bảng câu hỏi 2023 (Brouthers et al., 2016;
Goldfarb and Tucker, 2019).

Đối với doanh nghiệp tham gia thị trường nước ngoài, ngay cả việc áp
dụng website nền tảng cũng có thể có giá trị. Một website giảm chi phí
để khách hàng nước ngoài định vị doanh nghiệp, hỗ trợ giao tiếp bất đồng
bộ qua các múi giờ, và phát tín hiệu về tính hợp pháp cơ bản trong giao
dịch quốc tế. Phạm vi của những lợi ích này có giới hạn, một website tự
nó không tích hợp giao dịch, chuỗi cung ứng hay quyết định, và mức sinh
lời thực tế phụ thuộc vào việc doanh nghiệp có quy mô và thông lệ để
chuyển hoá khả năng hiển thị trực tuyến thành hoạt động kinh doanh xuất
khẩu hay không.

Hiện diện số dựa trên website do đó nên cho thấy mối liên hệ thuận chiều
trung bình với hiệu quả, nhưng không nhất thiết đồng nhất qua các đợt và
giai đoạn. Giá trị thông tin của một website có thể đã dịch chuyển qua
cửa sổ 2009--2023:

vào năm 2009, một website là một giao diện thị trường khác biệt; đến năm
2023 nó gần với một dấu hiệu thường lệ của hiện diện số cơ bản. Phần
thảo luận trong mục 5 xem xét nghiêm túc cách đọc "biến đại diện đã lỗi
thời" (proxy obsolescence) như một lựa chọn thay thế cho câu chuyện
thuần tuý về phụ thuộc giai đoạn.

Theo Verhoef et al. (2021), năng lực số có thể được định vị trên một thứ
bậc bốn tầng: Tầng 1, hiện diện số (website, email); Tầng 2, giao tiếp
số và thương mại điện tử cơ bản; Tầng 3, tích hợp quy trình số (thanh
toán điện tử, số hoá chuỗi cung ứng); Tầng 4, năng lực số động (ra quyết
định dựa trên dữ liệu, tích hợp AI). \(\tilde{\mathrm{DAI}}\) chính chỉ
neo ở Tầng 1. Một phần mở rộng kiểu Tầng 3
(\(\mathrm{DAI}_{\text{rich}}\)) được báo cáo trong mục 4.5 Panel B cho
đợt 2023, nơi các chỉ báo thanh toán điện tử trở nên sẵn có. Nhãn áp
dụng website nền tảng do đó theo sát những gì khái niệm có thể thực sự
nhận diện trên cửa sổ 2009--2023.

Vì \(\tilde{\mathrm{DAI}}\) được neo ở Tầng 1 (chỉ hiện diện website),
khái niệm này không tương đương với các thước đo áp dụng công nghệ số
cũng ghi nhận các chỉ báo Tầng 2 kích hoạt giao dịch như cường độ thanh
toán điện tử. Trong các bối cảnh nơi hạ tầng số Tầng 2--3 đã trưởng
thành và phổ cập, áp dụng công nghệ số cơ bản có thể tương tác với cường
độ xuất khẩu thông qua cơ chế khác. Cụ thể, nó vận hành như bổ trợ mở
rộng có điều kiện thay vì khuếch đại căng thẳng phối hợp, vì hệ sinh
thái xung quanh có thể hấp thụ xử lý giao dịch xuyên biên giới mà một
website đơn lẻ không thể.

\textbf{Giả thuyết H3.} Hiện diện số dựa trên website
(\(\tilde{\mathrm{DAI}}\)) có liên hệ thuận chiều với hiệu quả hoạt động
kinh doanh của doanh nghiệp ở Việt Nam, tính trung bình.

\hypertarget{24-giuxe1-trux1ecb-sux1ed1-phux1ee5-thuux1ed9c-giai-ux111oux1ea1n}{%
\subsubsection{2.4 Giá trị số phụ thuộc giai
đoạn}\label{24-giuxe1-trux1ecb-sux1ed1-phux1ee5-thuux1ed9c-giai-ux111oux1ea1n}}

Tuyên bố lý thuyết cốt lõi của nghiên cứu này là áp dụng công nghệ số
nền tảng phụ thuộc giai đoạn. Trong các giai đoạn đầu của quốc tế hoá,
công cụ số có thể tạo ra lợi ích tương đối trực tiếp bằng cách giúp
doanh nghiệp giao tiếp nhanh hơn, tiếp cận thị trường dễ hơn và quản lý
giao dịch hiệu quả hơn. Tuy nhiên, ở các giai đoạn sau, lợi ích của số
hoá có thể trở nên có điều kiện hơn vì doanh nghiệp đối mặt với các yêu
cầu phối hợp phức tạp hơn.

Trong điều kiện đó, giá trị của áp dụng công nghệ số nền tảng ngày càng
phụ thuộc vào việc công cụ số có được nhúng vào các thông lệ tổ chức
rộng hơn và phù hợp với quy mô xuất khẩu hay không (Banalieva and
Dhanaraj, 2019; Petricevic and Teece, 2019).

Lập luận này ngụ ý rằng số hoá có thể là con dao hai lưỡi trong một nền
kinh tế chuyển đổi. Nó có thể tạo ra lợi ích có thể quan sát, nhưng
những lợi ích đó không đồng đều và phụ thuộc vào thời điểm, quy mô và
năng lực bổ trợ. Trong một số giai đoạn, áp dụng công nghệ số có thể chủ
yếu vận hành như một yếu tố trực tiếp nâng cao hiệu quả. Trong các giai
đoạn khác, nó có thể suy yếu, biến mất hoặc trở nên có điều kiện với mức
độ quốc tế hoá của doanh nghiệp.

Kỳ vọng này đặc biệt phù hợp ở Việt Nam, nơi doanh nghiệp hoạt động
trong một môi trường chuyển đổi chứ không phải ổn định thể chế và năng
lực hoàn toàn. Cách diễn giải vòng đời do đó phù hợp hơn cách diễn giải
phụ trội đồng nhất.

Hai cân nhắc khác làm sắc nét hơn dự đoán. Thứ nhất, khi nhóm doanh
nghiệp xuất khẩu tập trung ở sản xuất cường độ xuất khẩu thấp, công cụ
số nền tảng chủ yếu thực hiện vai trò tiếp cận thị trường: chúng giúp
doanh nghiệp tìm khách hàng, truyền đạt thông tin về giá và sản phẩm, và
xử lý giao dịch đơn giản. Lợi ích năng suất biên từ vai trò này là dương
nhưng phân bổ rộng rãi qua phạm vi cường độ xuất khẩu.

Thứ hai, khi nhóm doanh nghiệp xuất khẩu dịch chuyển về phía các doanh
nghiệp hoạt động ở cường độ xuất khẩu cao hơn và tham gia phối hợp xuyên
biên giới chặt chẽ hơn, các công cụ số Tầng 1--2 cùng loại bắt đầu tương
tác với chi phí phối hợp biên của các thị trường nước ngoài bổ sung.
Tương tác này có thể có tính thay thế (công cụ số làm giảm chi phí phối
hợp và khuếch đại cổ tức năng suất) hoặc bổ sung với lợi suất giảm dần
(công cụ số ở cường độ cao bộc lộ sự thiếu tích hợp sâu hơn và khuếch
đại căng thẳng phối hợp). Đây về cơ bản là câu hỏi thực nghiệm mà bài
báo này xem như kiểm định cho H4.

\textbf{Giả thuyết H4 (khám phá).} Tính liên quan đến năng suất của áp
dụng công nghệ số nền tảng (\(\tilde{\mathrm{DAI}}\)) biến đổi qua các
giai đoạn quốc tế hoá và chuyển đổi thể chế. Bất kỳ điều tiết nào của
đường cong cường độ xuất khẩu bởi áp dụng công nghệ số nền tảng do đó
được kỳ vọng đặc trưng theo đợt khảo sát thay vì hiện diện đồng nhất
xuyên suốt các giai đoạn, với khả năng phát hiện trong mẫu mạnh nhất
được dự kiến vào năm 2023.

Năng lực công nghệ (\(\tilde{\mathrm{TCI}}\)) H2 H2 (điều tiết)

Quốc tế hoá → H1 (không đơn điệu) → Hiệu quả hoạt động doanh nghiệp
(\(\mathrm{FSTS}^{c}\), \((\mathrm{FSTS}^{c})^2\)) (ln năng suất lao
động) H4 (điều tiết, khám phá) H3

Áp dụng công nghệ số nền tảng (\(\tilde{\mathrm{DAI}}\))

Các biến kiểm soát: \(\ln(\mathrm{Emp})\), FirmAge, ForeignOwned, hiệu
ứng cố định ngành (sector FE) {[}+ hiệu ứng cố định đợt khảo sát (wave
FE){]}

\includegraphics{figures/p3_vietnam/figure_1_conceptual_model.png}

\emph{Hình 1.} Mô hình khái niệm. Biến độc lập (quốc tế hoá,
\(\mathrm{FSTS}^{c}\) và \((\mathrm{FSTS}^{c})^2\)) và biến phụ thuộc
(hiệu quả hoạt động kinh doanh của doanh nghiệp, ln năng suất lao động)
tạo nên trục IV--DV. Năng lực công nghệ (\(\tilde{\mathrm{TCI}}\)) và áp
dụng công nghệ số nền tảng (\(\tilde{\mathrm{DAI}}\)) đóng vai trò tác
động trực tiếp (H2, H3) và là các biến điều tiết của đường cong FSTS
(điều tiết H2, điều tiết khám phá H4). Các biến kiểm soát được đưa vào
theo dạng cộng. Hiệu ứng cố định đợt khảo sát chỉ áp dụng cho đặc tả
gộp.

\hypertarget{3-dux1eef-liux1ec7u-vuxe0-phux1b0ux1a1ng-phuxe1p}{%
\subsection{3. Dữ liệu và phương
pháp}\label{3-dux1eef-liux1ec7u-vuxe0-phux1b0ux1a1ng-phuxe1p}}

\hypertarget{31-cux1ea5u-truxfac-dux1eef-liux1ec7u}{%
\subsubsection{3.1 Cấu trúc dữ
liệu}\label{31-cux1ea5u-truxfac-dux1eef-liux1ec7u}}

Phân tích thực nghiệm sử dụng bằng chứng cấp doanh nghiệp đã được hài
hoà cho Việt Nam qua ba đợt của Khảo sát Doanh nghiệp của Ngân hàng Thế
giới: 2009, 2015 và 2023 (World Bank, 2010, 2016, 2024). Ước lượng mô
hình riêng theo từng đợt cho phép quan sát các mối quan hệ là ổn định
hay đặc trưng theo thời gian, trong khi ước lượng gộp xác định các tác
động trung bình trên toàn bộ giai đoạn rộng hơn.

Mẫu ước lượng hữu hiệu biến đổi giữa các đặc tả mô hình do các giá trị
thiếu trong các biến năng lực. Trong các mô hình theo đợt đầy đủ, các
mẫu sử dụng được là 989 quan sát cho năm 2009, 956 cho năm 2015, và
1.013 cho năm 2023. Mô hình gộp đầy đủ chứa 2.958 quan sát. Cấu trúc này
cung cấp đủ biến thiên để so sánh các tác động trực tiếp và các mẫu hình
có điều kiện giữa các giai đoạn.

\hypertarget{32-biux1ebfn-sux1ed1}{%
\subsubsection{3.2 Biến số}\label{32-biux1ebfn-sux1ed1}}

Hiệu quả hoạt động doanh nghiệp được đo bằng năng suất lao động dưới
dạng logarit (\(\ln(\mathrm{LP})\)). Quốc tế hoá được đo bằng cường độ
xuất khẩu trực tiếp (FSTS), trừ trung bình trong từng đợt
(\(\mathrm{FSTS}^{c}\)) và bình phương (\((\mathrm{FSTS}^{c})^2\)) để có
thể đưa đồng thời các số hạng tuyến tính và bậc hai nhằm kiểm định tính
phi tuyến.

Phân tích sử dụng hai khái niệm năng lực phân biệt. Chỉ số năng lực công
nghệ (\(\tilde{\mathrm{TCI}}\)) được diễn giải hẹp như thước đo năng lực
công nghệ nước ngoài và tiêu chuẩn. Nó gồm chứng nhận chất lượng quốc tế
và công nghệ được cấp phép nước ngoài, đại diện cho mức độ tiếp xúc của
doanh nghiệp với các tiêu chuẩn công nghệ bên ngoài, chứ không phải toàn
bộ trữ lượng năng lực hấp thụ kiểu Cohen--Levinthal (Lall, 1992; Cohen
and Levinthal, 1990). Chỉ số áp dụng công nghệ số
(\(\tilde{\mathrm{DAI}}\)) được diễn giải hẹp như hiện diện số dựa trên
website (áp dụng website nền tảng), không phải là thước đo tích hợp số ở
cấp giao dịch hay chuyển đổi số (Bharadwaj et al., 2013; Verhoef et al.,
2021; Nambisan et al., 2019). Mỗi chỉ số hợp thành được chuẩn hoá theo z
trong từng đợt để các hệ số báo cáo có thể so sánh theo đơn vị độ lệch
chuẩn. Ký hiệu dấu ngã \(\tilde{\mathrm{TCI}}\) và
\(\tilde{\mathrm{DAI}}\) đánh dấu biến thể đã chuẩn hoá. Sự phân tách
này có chủ ý: nghiên cứu quan tâm liệu hai miền có thể hiện vai trò thực
nghiệm khác nhau hay không, đồng thời giữ chặt các nhãn khái niệm đúng
với những gì các chỉ báo WBES nền thực sự đo lường.

Cấu trúc cấp chỉ báo như sau. Biến kết quả là \(\ln(\mathrm{LP})\) =
ln(d2 / l1), trong đó d2 là tổng doanh thu hàng năm và l1 là số lao động
chính thức toàn thời gian. Quốc tế hoá là FSTS = d3c / 100, trừ trung
bình trong từng đợt (\(\mathrm{FSTS}^{c}\)) và bình phương
(\((\mathrm{FSTS}^{c})^2\)). Chỉ số năng lực công nghệ chính
(\(\tilde{\mathrm{TCI}}\)) là trung bình đã chuẩn hoá theo z trong từng
đợt của b8 (chứng nhận chất lượng được công nhận quốc tế) và e6 (công
nghệ được cấp phép nước ngoài), mỗi chỉ báo được mã hoá lại từ dạng 1/2
sang dạng nhị phân 1/0 của WBES.

Chỉ số áp dụng công nghệ số chính (\(\tilde{\mathrm{DAI}}\)) là chỉ báo
hiện diện website c22b đã chuẩn hoá theo z trong từng đợt, mã hoá lại từ
1/2 sang 1/0. Theo đặc tả chính đã sửa đổi này, không có chỉ báo nào
được dùng chung giữa các chỉ số hợp thành TCI và DAI; e6 chỉ thuộc về
khái niệm năng lực, trong khi c22b đơn lẻ đóng vai trò là biến đại diện
hài hoà giữa các đợt cho hiện diện số cơ bản. Các chỉ báo số cấp giao
dịch có thể so sánh giữa các đợt không có trong bộ công cụ 2009 và 2015,
do đó \(\tilde{\mathrm{DAI}}\) nên được diễn giải như một thước đo áp
dụng công nghệ số Tầng 1--2 chứ không phải một khái niệm năng lực số
tích hợp đầy đủ (Bharadwaj et al., 2013; Verhoef et al., 2021).

Hai chỉ số hợp thành được làm giàu được sử dụng trong các bảng kiểm tra
độ vững ở mục 4.5 khi tính sẵn có của các chỉ báo cho phép.
\(\mathrm{TCI}_{\text{full}}\) bổ sung h1 (đã giới thiệu sản phẩm mới
hoặc cải tiến đáng kể) và h8 (chỉ báo chi tiêu R\&D) vào các chỉ báo TCI
trong năm 2015 và 2023, nơi h1 và h8 hiện diện.
\(\mathrm{DAI}_{\text{rich}}\), chỉ được xây dựng cho đợt 2023, mở rộng
chỉ báo hiện diện website với k33 (tỉ trọng doanh thu nhận qua thanh
toán điện tử) và k38 (tỉ trọng thanh toán nhà cung cấp thực hiện qua
thanh toán điện tử), ở cả dạng liên tục (k33 / 100, k38 / 100) và dạng
nhị phân (k33 \textgreater{} 0, k38 \textgreater{} 0).
\(\mathrm{DAI}_{\text{rich}}\) do đó dịch chuyển từ hiện diện số cơ bản
về phía một khái niệm áp dụng công nghệ số kích hoạt giao dịch, nhưng
được xem là một kiểm tra độ vững về chiều sâu đo lường hơn là đặc tả
chính vì k33 và k38 không có trong các đợt 2009 và 2015.

\textbf{Ghi chú về tính so sánh giữa các tầng khái niệm.} Khái niệm
\(\tilde{\mathrm{DAI}}\) trong bài báo này là một chỉ báo chỉ ở Tầng 1
(nhị phân website c22b). Một chỉ số hợp thành phong phú hơn Tầng 1+2,
kết hợp hiện diện website với cường độ thanh toán điện tử (k33/k38), về
nguyên tắc sẽ ghi nhận các cơ chế kích hoạt giao dịch mở rộng theo các
yêu cầu phối hợp xuyên biên giới, nhưng k33 và k38 không có trong các
đợt 2009 và 2015 của WBES Việt Nam, loại trừ một đặc tả Tầng 1+2 giữa
các đợt. Tương tác Tầng 1 đơn lẻ âm được ghi nhận ở đây
(\(\mathrm{FSTS}^{c}\) × \(\tilde{\mathrm{DAI}}\) = −0,912, p = .043) do
đó được diễn giải như tính lỗi thời của tầng khái niệm: hiện diện
website Tầng 1 đã trở thành một chứng chỉ ngưỡng tối thiểu trong môi
trường số đang trưởng thành của Việt Nam và không còn phân biệt năng lực
phối hợp xuyên biên giới của doanh nghiệp ở cường độ xuất khẩu cao. Chỉ
số hợp thành kiểm tra độ vững \(\mathrm{DAI}_{\text{rich}}\) có sẵn
trong đợt 2023 phần nào lấp khoảng trống này cho phân tích nhạy cảm
trong đợt, nhưng tính so sánh giữa các đợt giới hạn việc sử dụng nó như
thước đo chính.

Các biến kiểm soát theo chuẩn. Quy mô doanh nghiệp \(\ln(\mathrm{Emp})\)
= ln(l1). Tuổi doanh nghiệp FirmAge = năm khảo sát trừ b5 (năm thành
lập). Sở hữu nước ngoài ForeignOwned = 1 nếu b2b (tỉ lệ vốn cổ phần do
cá nhân hoặc doanh nghiệp nước ngoài tư nhân sở hữu) \textgreater{} 0.
Hiệu ứng cố định ngành sử dụng chữ số đầu tiên của a4b (mã ISIC rộng)
cho các đợt 2009 và 2015 và chữ số đầu tiên của a4a cho đợt 2023, nơi
a4b không có trong bản công bố công khai. Các đặc tả gộp bổ sung hiệu
ứng cố định đợt khảo sát để hấp thụ các khác biệt giai đoạn rộng trong
mức năng suất nền.

Ghi chú về xử lý mã thiếu. Các mã công cụ WBES ghi nhận các phản hồi
không biết và từ chối là -9. Chúng tôi xem -9 là thiếu trước khi xây
dựng bất kỳ chỉ số hợp thành nào và áp dụng loại bỏ theo dòng (listwise
deletion) trên tập biến trọng tâm (\(\ln(\mathrm{LP})\),
\(\ln(\mathrm{Emp})\), FirmAge, ForeignOwned, FSTS,
\(\mathrm{TCI}_{\text{thin}}\), \(\mathrm{DAI}_{\text{thin}}\),
sector1). Các mẫu phân tích thu được là 989, 956 và 1.013 quan sát cho
các năm 2009, 2015 và 2023 tương ứng; N gộp là 2.958.

\hypertarget{33-truxecnh-tux1ef1-muxf4-huxecnh}{%
\subsubsection{3.3 Trình tự mô
hình}\label{33-truxecnh-tux1ef1-muxf4-huxecnh}}

Chiến lược thực nghiệm theo một chuỗi lồng nhau các mô hình OLS được ước
lượng với sai số chuẩn vững đối với phương sai thay đổi HC1 (MacKinnon
\& White, 1985). Đặc tả đầy đủ quan tâm là

\[
\begin{aligned}
\ln(\mathrm{LP}_i) =\;& \alpha + \beta_1\,\mathrm{FSTS}^{c}_i + \beta_2\,(\mathrm{FSTS}^{c}_i)^2 + \beta_3\,\tilde{\mathrm{TCI}}_i + \beta_4\,\tilde{\mathrm{DAI}}_i \\
& + \beta_5\,(\mathrm{FSTS}^{c}_i \times \tilde{\mathrm{DAI}}_i) + \beta_6\,((\mathrm{FSTS}^{c}_i)^2 \times \tilde{\mathrm{DAI}}_i) \\
& + \boldsymbol{\gamma}^{\top}\!\mathbf{x}_i + \delta_s + \lambda_w + \varepsilon_i,
\end{aligned}
\]

trong đó
\(\mathrm{FSTS}^{c}_i = \mathrm{FSTS}_i - \overline{\mathrm{FSTS}}_w\)
là cường độ xuất khẩu được trung tâm hoá theo đợt, \(\mathbf{x}_i\) tập
hợp các biến kiểm soát cấp doanh nghiệp
\(\{\mathrm{lnEmp}_i, \mathrm{FirmAge}_i, \mathrm{ForeignOwned}_i\}\),
\(\delta_s\) là hiệu ứng cố định ngành ISIC một chữ số, và \(\lambda_w\)
là hiệu ứng cố định đợt khảo sát chỉ áp dụng trong đặc tả gộp. Điểm
ngoặt ngụ ý của chữ U ngược là
\(\mathrm{FSTS}^{*} = -\hat\beta_1/(2\hat\beta_2) + \overline{\mathrm{FSTS}}_w\)
khi \(\hat\beta_2 < 0\) (Haans, Pieters, \& He, 2016).

Các đặc tả lồng nhau được ước lượng để tách các mối liên hệ trực tiếp
khỏi các mối liên hệ có điều kiện:

\begin{itemize}
\tightlist
\item
  \textbf{M0} (chỉ có các biến kiểm soát):
  \(\ln(\mathrm{LP}_i) = \alpha + \boldsymbol{\gamma}^{\top}\!\mathbf{x}_i + \delta_s + \lambda_w + \varepsilon_i\);
\item
  \textbf{M1} (FSTS tuyến tính): bổ sung
  \(\beta_1\,\mathrm{FSTS}^{c}_i\);
\item
  \textbf{M2} (chữ U ngược): bổ sung
  \(\beta_2\,(\mathrm{FSTS}^{c}_i)^2\);
\item
  \textbf{M7} (năng lực trực tiếp kép): M2 cộng
  \(\beta_3\,\tilde{\mathrm{TCI}}_i + \beta_4\,\tilde{\mathrm{DAI}}_i\);
\item
  \textbf{M8} (điều tiết DAI đầy đủ): M7 cộng hai số hạng tương tác DAI
  ở trên.
\end{itemize}

Các khác biệt hệ số giữa các đợt được đánh giá thông qua kiểm định z
Paternoster et al. (1998),
\(z = (\hat\beta_A - \hat\beta_B)/\sqrt{\mathrm{SE}(\hat\beta_A)^2 + \mathrm{SE}(\hat\beta_B)^2}\),
với các giá trị \(p\) hai phía từ phân phối chuẩn tiêu chuẩn. Kiểm định
U của Lind và Mehlum (2010) được áp dụng trên khoảng \([0, 1]\) của FSTS
để xác nhận chính thức hình chữ U ngược.

Thiết kế này tách ba câu hỏi phân tích. Thứ nhất, mối quan hệ I--P có
phi tuyến hay không? Thứ hai, \(\tilde{\mathrm{TCI}}\) và
\(\tilde{\mathrm{DAI}}\) có liên hệ trực tiếp với hiệu quả hay không?
Thứ ba, vai trò của áp dụng công nghệ số nền tảng có trở nên có điều
kiện hơn khi cường độ xuất khẩu tăng hay không? Xuyên suốt, các kết quả
được mô tả như các mối liên hệ thay vì tác động, nhất quán với giới hạn
suy luận của dữ liệu cắt ngang lặp lại (Antonakis et al., 2010;
Wooldridge, 2010).

\hypertarget{34-tuxe1i-lux1eb7p-vuxe0-khux1ea3-nux103ng-tuxe1i-tux1ea1o}{%
\subsubsection{3.4 Tái lặp và khả năng tái
tạo}\label{34-tuxe1i-lux1eb7p-vuxe0-khux1ea3-nux103ng-tuxe1i-tux1ea1o}}

Toàn bộ quy trình được triển khai dưới dạng một sơ đồ Stata 10 bước được
phân phối kèm bản thảo. Các bước xây dựng (01--04) làm sạch từng đợt
WBES, hài hoà tập biến trọng tâm và nối ba đợt vào một tệp gộp với việc
trung tâm hoá trong từng đợt và chuẩn hoá theo z được áp dụng lại. Các
bước ước lượng (05--09) bao gồm chuỗi lồng nhau M0--M8, kiểm tra điểm
ngoặt Lind--Mehlum, các kiểm định chọn lựa Heckman thủ công, các kiểm
định z giữa các đợt theo Paternoster et al. (1998), và các bảng kiểm tra
độ vững được mô tả trong mục 4.5. Bước xuất (10) ghi các bảng và Hình 2
hướng tới bản thảo trực tiếp từ các ước lượng đã lưu.

Việc chạy lại quy trình từ một bản nhân bản mới sẽ tái tạo mọi hệ số
được báo cáo bên dưới; văn bản bản thảo chứ không phải kết quả tệp
do-file là đối tượng điều chỉnh khi việc chạy lại trôi khỏi văn bản.
Xuyên suốt, chúng tôi tuân theo các khuyến nghị thực hành tốt nhất hiện
hành về chuẩn bị dữ liệu, tính minh bạch và diễn giải tích luỹ của bằng
chứng liên hệ (Aguinis et al., 2021; Shaver, 2020).

\hypertarget{4-kux1ebft-quux1ea3}{%
\subsection{4. Kết quả}\label{4-kux1ebft-quux1ea3}}

Bảng 1 báo cáo thống kê tóm tắt của mẫu phân tích theo từng đợt. Có ba
mẫu hình đáng chú ý trước khi tiến hành phân tích suy luận. Tỉ lệ doanh
nghiệp báo cáo cường độ xuất khẩu trực tiếp dương giảm từ 28,4 \% năm
2009 xuống 20,7 \% năm 2015 và 18,8 \% năm 2023. Mô thức này phản ánh sự
tái cân bằng của nhóm doanh nghiệp xuất khẩu Việt Nam từ sản xuất thâm
dụng lao động sang các doanh nghiệp dịch vụ và doanh nghiệp chuỗi cung
ứng liên kết FDI trong cửa sổ quan sát. Mức trung bình trong đợt của áp
dụng công nghệ số cơ bản (chỉ báo website c22b) tăng từ 0,425 năm 2009
lên 0,483 năm 2015 và 0,498 năm 2023, cho thấy sự lan toả rộng rãi của
hiện diện số cơ bản trong dân doanh nghiệp Việt Nam qua 14 năm. Mức
trung bình của năng suất lao động logarit tăng đơn điệu
(19,41/20,04/20,55), nhất quán với hội tụ năng suất rộng hơn của Việt
Nam trong giai đoạn.

\hypertarget{41-phuxe1t-hiux1ec7n-theo-tux1eebng-ux111ux1ee3t}{%
\subsubsection{4.1 Phát hiện theo từng
đợt}\label{41-phuxe1t-hiux1ec7n-theo-tux1eebng-ux111ux1ee3t}}

Đợt 2009 cho thấy một mối quan hệ quốc tế hoá--hiệu quả phi tuyến rõ rệt
cùng với các tác động năng lực trực tiếp và áp dụng công nghệ số mạnh.
Đặc tả chữ U ngược (M2) cho ra số hạng tuyến tính dương (β = 1,045, p =
.015) và số hạng bậc hai âm (β = -1,774, p = .009); kiểm định
Lind--Mehlum bác bỏ tính đơn điệu tại p = .006. Trong đặc tả trực tiếp
kép (M7), cả \(\tilde{\mathrm{TCI}}\) (β = 0,215, p \textless{} .001) và
\(\tilde{\mathrm{DAI}}\) (β = 0,175, p \textless{} .001) đều dương và có
ý nghĩa cao. Điều tiết của TCI có thể phân biệt thống kê với không (M3
joint p = .040; \(\mathrm{FSTS}^{c}\) × \(\tilde{\mathrm{TCI}}\) =
-0,579, p = .087), nhưng điều tiết của DAI thì không (M4 joint p = .825;
M8 joint đầy đủ p = .700). Về mặt thực chất, cả hai chiều năng lực trong
2009 vận hành chủ yếu như các yếu tố dịch chuyển mức (level shift) trực
tiếp; chi phí phối hợp biên làm uốn cong đường I--P có liên hệ với năng
lực công nghệ chứ không phải với hiện diện số cơ bản.

Bảng 1: Thống kê tóm tắt mẫu phân tích theo từng đợt.

\begin{longtable}[]{@{}lllll@{}}
\toprule\noalign{}
Biến số & 2009 (N=989) & 2015 (N=956) & 2023 (N=1.013) & Gộp
(N=2.958) \\
\midrule\noalign{}
\endhead
\bottomrule\noalign{}
\endlastfoot
Ln(Năng suất lao động) & 19,412 (1,307) & 20,042 (1,460) & 20,549
(1,474) & 20,005 (1,491) \\
FSTS (cường độ xuất khẩu) & 0,168 (0,337) & 0,119 (0,283) & 0,131
(0,311) & 0,139 (0,312) \\
Tỉ lệ doanh nghiệp xuất khẩu & 0,284 & 0,207 & 0,188 & 0,226 \\
\(\tilde{\mathrm{TCI}}\) (trung bình) & 0,169 (0,305) & 0,142 (0,295) &
0,146 (0,276) & 0,152 (0,292) \\
\(\tilde{\mathrm{DAI}}\) (trung bình) & 0,425 (0,495) & 0,483 (0,500) &
0,498 (0,500) & 0,469 (0,499) \\
Quy mô doanh nghiệp (ln số lao động) & 4,067 (1,493) & 3,629 (1,476) &
3,578 (1,539) & 3,758 (1,519) \\
Tuổi doanh nghiệp (năm) & 11,900 (11,300) & 12,800 (9,600) & 14,100
(7,900) & 12,900 (9,700) \\
Sở hữu nước ngoài (tỉ lệ) & 0,142 & 0,090 & 0,125 & 0,119 \\
\end{longtable}

\emph{Ghi chú.} Báo cáo trung bình (SD). \(\tilde{\mathrm{TCI}}\) và
\(\tilde{\mathrm{DAI}}\) là các chỉ số hợp thành cấu thành đã chuẩn hoá
theo z. FSTS = tỉ lệ doanh thu nước ngoài trên tổng doanh thu. Cỡ mẫu:
2009 N=989; 2015 N=956; 2023 N=1.013; gộp N=2.958.

Bảng 2: Tổng hợp hệ số trọng tâm theo từng đợt và mẫu gộp (M2 chữ U
ngược, M7 trực tiếp kép, M8 với điều tiết DAI).

\begin{longtable}[]{@{}lllll@{}}
\toprule\noalign{}
Số hạng & 2009 (N=989) & 2015 (N=956) & 2023 (N=1.013) & Gộp
(N=2.958) \\
\midrule\noalign{}
\endhead
\bottomrule\noalign{}
\endlastfoot
\textbf{M2. Chữ U ngược trên FSTS} & & & & \\
\(\mathrm{FSTS}^{c}\) (tuyến tính) & +1,045* (p=.015) & +1,159* (p=.029)
& +0,962* (p=.039) & +0,984*** (p\textless.001) \\
\((\mathrm{FSTS}^{c})^2\) (bậc hai) & -1,823** (p=.005) & -2,115**
(p=.004) & -1,686** (p=.008) & -1,909*** (p\textless.001) \\
Lind--Mehlum p & \textless.05 & .009 & .013 & \textless.001 \\
Điểm ngoặt (FSTS) & 0,287 & 0,274 & 0,285 & 0,397 \\
\textbf{M7. Năng lực trực tiếp kép} & & & & \\
\(\tilde{\mathrm{TCI}}\) & +0,215*** (p\textless.001) & +0,128**
(p=.010) & +0,123** (p=.006) & +0,179*** (p\textless.001) \\
\(\tilde{\mathrm{DAI}}\) & +0,175*** (p\textless.001) & -0,044 (p=.377)
& +0,095* (p=.038) & +0,078** (p=.004) \\
\textbf{M8. Với điều tiết FSTS × DAI (M8 gộp: \(\mathrm{FSTS}^{c}\)
β=0,845, p=.006; \((\mathrm{FSTS}^{c})^2\) β=-1,650, p\textless.001)} &
& & & \\
\(\mathrm{FSTS}^{c}\) × \(\tilde{\mathrm{DAI}}\) & n.s. & n.s. & -0,912*
(p=.043) & (lẫn lộn giữa các đợt) \\
\((\mathrm{FSTS}^{c})^2\) × \(\tilde{\mathrm{DAI}}\) & n.s. & n.s. &
n.s. & n.s. \\
Kiểm định liên hợp M8 (độ cong + điều tiết) & p=.700 & n.s. & p=.062 † &
--- \\
\end{longtable}

\emph{Ghi chú.} Hệ số từ OLS với sai số chuẩn vững đối với phương sai
thay đổi HC1. Hiệu ứng cố định ngành và các biến kiểm soát cấp doanh
nghiệp chuẩn (\(\ln(\mathrm{Emp})\), FirmAge, ForeignOwned) được bao gồm
xuyên suốt. Ý nghĩa: *** p \textless{} .001; ** p \textless{} .01; * p
\textless{} .05; † p \textless{} .10. M2 = chỉ FSTS bậc hai; M7 = M2 +
\(\tilde{\mathrm{TCI}}\) + \(\tilde{\mathrm{DAI}}\) (trực tiếp kép); M8
= M7 + \(\mathrm{FSTS}^{c}\) × \(\tilde{\mathrm{DAI}}\) +
\((\mathrm{FSTS}^{c})^2\) × \(\tilde{\mathrm{DAI}}\). Đợt 2023 là đợt
duy nhất trong đó kênh điều tiết DAI của M8 có thể phát hiện được ở các
ngưỡng thông thường, nhất quán với cách diễn giải tính lỗi thời của biến
đại diện Tầng 1 được phát triển trong Mục 5.

Đợt 2015 cho thấy độ cong rõ ràng nhưng kênh số yếu nhất. M2 cho ra
\(\mathrm{FSTS}^{c}\) β = 1,159 (p = .029) và \((\mathrm{FSTS}^{c})^2\)
β = -2,115 (p = .004), với Lind--Mehlum p = .009.
\(\tilde{\mathrm{TCI}}\) vẫn giữ một mối liên hệ trực tiếp thuận chiều
(β = 0,128, p = .010) nhưng chỉ ở khoảng 60 \% độ lớn của năm 2009.
\(\tilde{\mathrm{DAI}}\) mất hoàn toàn ý nghĩa trực tiếp (β = -0,044, p
= .377). Điều tiết TCI không có ý nghĩa trong đợt này (M3 joint p =
.713), và điều tiết DAI cao nhất chỉ ở mức cận ngưỡng (M4 joint p =
.125; M8 joint p = .093). Đọc như đặc trưng giai đoạn, năm 2015 trông
giống đợt trong đó độ cong I--P sắc nét bất thường trong khi kênh số nén
hoàn toàn. Mẫu hình này nhất quán với câu chuyện đường cong J về năng
suất, trong đó doanh nghiệp đầu tư vào công cụ số cơ bản trước khi các
điều chỉnh tổ chức bổ trợ tạo ra lợi ích năng suất đo lường được
(Brynjolfsson et al., 2021).

Đợt 2023 là nơi tín hiệu điều tiết số nổi lên rõ nét nhất.

M2 lại chỉ ra một chữ U ngược rõ ràng (\(\mathrm{FSTS}^{c}\) β = 0,962,
p = .039; \((\mathrm{FSTS}^{c})^2\) β = -1,686, p = .008; Lind--Mehlum p
= .013).

Trong đặc tả M7 trực tiếp kép, cả hai chiều năng lực đều dương và có ý
nghĩa (\(\tilde{\mathrm{TCI}}\) β = 0,123, p = .006;
\(\tilde{\mathrm{DAI}}\) β = 0,095, p = .038). Khi các số hạng tương tác
DAI được bổ sung vào M8, tương tác tuyến tính âm và có ý nghĩa riêng lẻ
(\(\mathrm{FSTS}^{c}\) × \(\tilde{\mathrm{DAI}}\) = -0,912, p = .043),
tương tác bậc hai dương và cận ngưỡng (\((\mathrm{FSTS}^{c})^2\) ×
\(\tilde{\mathrm{DAI}}\) = 1,043, p = .099), và kiểm định liên hợp ở mức
cận ý nghĩa (M4 joint p = .102; M8 joint p = .062). Cách đọc thực chất
là, đến năm 2023, áp dụng công nghệ số cơ bản trở nên có điều kiện hơn
với cường độ xuất khẩu:

khi các doanh nghiệp xuất khẩu vượt qua các mức FSTS vừa phải, đóng góp
năng suất của hiện diện số cơ bản suy yếu và có thể khuếch đại thay vì
giảm bớt gánh nặng phối hợp. Phát hiện này nhất quán với quan điểm rằng
chỉ áp dụng ở cấp website là không đủ khi năng lực số kích hoạt giao
dịch và tích hợp quy trình sâu hơn vẫn đang trưởng thành.

Một cách diễn giải bổ sung là tính lỗi thời của biến đại diện. Cách này
nhất quán với kết quả 2SLS không có ý nghĩa cho DAI được công cụ hoá
(β=0,018, p=.942; Mục 4.5 Panel K). Đến năm 2023, hiện diện website
(Tầng 1, c22b) đã trở thành chứng chỉ ngưỡng tối thiểu chứ không phải
dấu hiệu phân biệt năng lực trong môi trường số đang trưởng thành của
Việt Nam. Các doanh nghiệp có và không có website không còn được phân
biệt một cách có ý nghĩa về năng lực phối hợp xuyên biên giới ở tầng DAI
mà c22b ghi nhận. Cách diễn giải này dự đoán tương tác DAI × FSTS âm mà
không yêu cầu rằng "áp dụng công nghệ số là xấu". Thay vào đó, công cụ
đã mất sức mạnh phân biệt khi áp dụng Tầng 1 lan toả đến mức gần phổ
quát (49,8 \% năm 2023 so với 42,5 \% năm 2009). Các phân tích tương lai
sử dụng chỉ số hợp thành \(\mathrm{DAI}_{\text{rich}}\) (Tầng 1+2; Panel
B, chỉ 2023) sẽ cung cấp kiểm định trong đợt về việc liệu dấu có dịch
chuyển khi các chỉ báo cường độ thanh toán điện tử được bổ sung vào công
cụ Việt Nam hay không.

Tóm lại, các kết quả theo từng đợt vẽ ra hai mối liên hệ đặc trưng theo
đợt nhất quán với phụ thuộc giai đoạn. Năng lực công nghệ nước ngoài /
tiêu chuẩn dương trong cả ba đợt với độ lớn suy giảm khiêm tốn
(\(\tilde{\mathrm{TCI}}\) = 0,215; 0,128; 0,123) nhưng với điều tiết
sống sót trong các bảng 2009, 2023 và gộp. Hiện diện số dựa trên website
theo một quỹ đạo không đơn điệu, mạnh trong năm 2009 (β = 0,175 ***),
không có ý nghĩa năm 2015 (β = -0,044 n.s.), và tái xuất hiện năm 2023
(β = 0,095 *), và kênh điều tiết của nó chỉ vật chất hoá vào năm 2023.
Các kiểm định z giữa các đợt theo Paternoster được báo cáo trong mục 4.5
xác nhận rằng mức giảm của DAI từ 2009 đến 2015 (z = 3,353, p
\textless{} .001) và sự phục hồi từ 2015 đến 2023 (z = -2,051, p = .040)
có thể phân biệt thống kê.

Một kiểm định tương tác chính thức gộp đợt × biến trọng tâm (báo cáo
dưới dạng Panel I trong mục 4.5) xác nhận rằng chỉ các dịch chuyển trực
tiếp DAI là có thể phân biệt giữa các đợt; các khác biệt độ cong FSTS và
điều tiết FSTS × DAI giữa các đợt không thể tách biệt thống kê trong đặc
tả gộp bão hoà. Chúng tôi do đó đọc mẫu hình như các mối liên hệ đặc
trưng theo đợt nhất quán với phụ thuộc giai đoạn, chứ không phải như một
dịch chuyển cấu trúc được nhận diện đầy đủ giữa các đợt.

Mẫu hình theo từng đợt mang một cách đọc thể chế. Đợt 2009 ghi nhận thời
kỳ ngay sau khi gia nhập WTO. Khi đó, doanh nghiệp xuất khẩu biên vẫn
nằm trong vùng chi phí nhập cuộc của đường I--P, trữ lượng năng lực là
nguồn lực khan hiếm, và các công cụ số nền tảng (ngay cả ở lớp chỉ riêng
website) tạo ra lợi ích trực tiếp vì lựa chọn thay thế là xử lý giao
dịch bằng giấy. Đợt 2015 ghi nhận giai đoạn chuyển đổi trong đó độ cong
I--P sắc nét bất thường (\((\mathrm{FSTS}^{c})^2\) = -2,115, p = .004
trong M2) nhưng kênh số nén hoàn toàn. \(\tilde{\mathrm{DAI}}\) mất ý
nghĩa trực tiếp và không có điều tiết liên hợp. Phát hiện này gợi ý rằng
các khác biệt trong thành phần doanh nghiệp xuất khẩu giữa các đợt phản
ánh ở việc lợi ích năng suất đến chủ yếu từ quy mô và năng lực công nghệ
nước ngoài/tiêu chuẩn chứ không phải từ áp dụng công nghệ số nền tảng
trong đợt này. Đợt 2023 ghi nhận sự tái xuất hiện của kênh số như biến
điều tiết chứ không phải khoản phụ trội trực tiếp đồng nhất. Hạ tầng số
sau NDTP khiến áp dụng công nghệ số nền tảng tương tác với cường độ xuất
khẩu. Tương tác \(\mathrm{FSTS}^{c}\) × \(\tilde{\mathrm{DAI}}\) âm cho
thấy kênh có điều kiện này ràng buộc chủ yếu ở cường độ xuất khẩu cao
hơn, chứ không đồng nhất trên phạm vi cường độ xuất khẩu.

\hypertarget{42-phuxe1t-hiux1ec7n-gux1ed9p}{%
\subsubsection{4.2 Phát hiện gộp}\label{42-phuxe1t-hiux1ec7n-gux1ed9p}}

Các ước lượng gộp xác nhận rằng mối quan hệ quốc tế hoá--hiệu quả là phi
tuyến tính trung bình. Trong đặc tả M2 gộp, số hạng tuyến tính
\(\mathrm{FSTS}^{c}\) dương (β = 0,984, p \textless{} .001) và số hạng
bậc hai \((\mathrm{FSTS}^{c})^2\) âm (β = -1,909, p \textless{} .001);
kiểm định Lind--Mehlum bác bỏ tính đơn điệu tại p \textless{} .001 với
điểm ngoặt ước lượng tại 39,7 \% của cường độ xuất khẩu trực tiếp. Độ
cong vẫn duy trì trong đặc tả M8 đầy đủ (\(\mathrm{FSTS}^{c}\) β =
0,845, p = .006; \((\mathrm{FSTS}^{c})^2\) β = -1,650, p \textless{}
.001), ủng hộ H1 và đồng nhất với bằng chứng phân tích tổng hợp về hình
dạng phi tuyến của lợi suất quốc tế hoá ở các doanh nghiệp thị trường
mới nổi (Marano et al., 2016).

Bằng chứng gộp cho thấy cả năng lực công nghệ và áp dụng công nghệ số cơ
bản đều có liên hệ thuận chiều với hiệu quả hoạt động kinh doanh của
doanh nghiệp tính trung bình. Trong đặc tả M7 trực tiếp kép gộp,
\(\tilde{\mathrm{TCI}}\) dương và có ý nghĩa cao (β = 0,179, p
\textless{} .001), và \(\tilde{\mathrm{DAI}}\) dương và có ý nghĩa (β =
0,078, p = .004).

Trong đặc tả M8 đầy đủ, hệ số \(\tilde{\mathrm{TCI}}\) về cơ bản không
thay đổi (β = 0,184, p \textless{} .001), trong khi hệ số trực tiếp
\(\tilde{\mathrm{DAI}}\) trở nên không thể phân biệt thống kê với không
(β = 0,032, p = .537) khi các số hạng tương tác được đưa vào. Tính nhạy
cảm này nhất quán với cách diễn giải rằng \(\tilde{\mathrm{DAI}}\) kết
hợp một tác động mức dương với một tương tác âm với
\(\mathrm{FSTS}^{c}\) ở cường độ xuất khẩu cao hơn, đúng mẫu hình được
ghi nhận trong mục 4.1 cho đợt 2023.

Các số hạng tương tác gộp liên quan đến \(\tilde{\mathrm{DAI}}\) mang
tín hiệu liên hợp cận ngưỡng. Tương tác tuyến tính âm nhưng không có ý
nghĩa riêng lẻ (\(\mathrm{FSTS}^{c}\) × \(\tilde{\mathrm{DAI}}\) =
-0,448, p = .116). Tương tác bậc hai dương nhưng không có ý nghĩa
(\((\mathrm{FSTS}^{c})^2\) × \(\tilde{\mathrm{DAI}}\) = 0,460, p =
.276). Kiểm định Wald liên hợp ở mức cận ý nghĩa (M8 joint p = .083).
Tín hiệu gộp này được chi phối chủ yếu bởi đợt 2023. Cụ thể, kiểm định
điều tiết liên hợp M4 trên DAI không có ý nghĩa trong năm 2009 (p =
.825) và 2015 (p = .125), và chỉ đạt mức cận ý nghĩa vào năm 2023 (M4
joint p = .102; M8 joint p = .062). Các ước lượng gộp do đó đánh giá
thấp thời điểm của kênh điều tiết DAI: kênh tập trung trong đợt gần đây
nhất chứ không phân bổ đồng đều xuyên cửa sổ 2009--2023.

Điều tiết năng lực công nghệ, ngược lại, được phân bổ đồng đều hơn.

Kiểm định liên hợp M3 trên \(\mathrm{FSTS}^{c}\) ×
\(\tilde{\mathrm{TCI}}\) và \((\mathrm{FSTS}^{c})^2\) ×
\(\tilde{\mathrm{TCI}}\) có thể phân biệt thống kê với không trong ba
trong bốn bảng (2009 p = .040, 2023 p = .027, gộp p = .003), và chỉ
không có ý nghĩa vào năm 2015 (p = .713). Trong mẫu gộp, tương tác tuyến
tính âm (\(\mathrm{FSTS}^{c}\) × \(\tilde{\mathrm{TCI}}\) = -0,587, p =
.003) và bậc hai dương (\((\mathrm{FSTS}^{c})^2\) ×
\(\tilde{\mathrm{TCI}}\) = 0,640, p = .031). Phát hiện này cho thấy chữ
U ngược phẳng lại đối với các doanh nghiệp có năng lực cao, chứ không
dịch chuyển về mức. Mô thức nhất quán với cách đọc năng lực hấp thụ: các
doanh nghiệp với trữ lượng năng lực sâu hơn rút lợi ích năng suất trên
phạm vi cường độ xuất khẩu rộng hơn so với các đồng nghiệp có năng lực
thấp hơn (Cohen and Levinthal, 1990; Lall, 1992).

Nếu phân tích dừng lại ở ước lượng gộp, người ta có thể kết luận rằng số
hoá cung cấp một khoản phụ trội hiệu quả tích cực rộng rãi nhưng đơn
giản về cấu trúc. Bằng chứng theo từng đợt cho thấy kết luận này sẽ
không đầy đủ.

Mức trung bình gộp dương cùng tồn tại với tính không đồng nhất theo thời
gian đáng kể, và vai trò có điều kiện của áp dụng công nghệ số chỉ nổi
lên rõ nét hơn trong đợt sau.

\includegraphics{figures/p3_vietnam/figure_2a.png}

\emph{Hình 2a.} Dự đoán ln(năng suất lao động) như một hàm của FSTS cho
đợt 2009 (M2). Vùng tô = khoảng tin cậy 95 \% (95 \% CI). Điểm ngoặt ≈
46 \% FSTS (Lind-Mehlum p = .006).

\includegraphics{figures/p3_vietnam/figure_2b.png}

\emph{Hình 2b.} Dự đoán ln(năng suất lao động) như một hàm của FSTS cho
đợt 2015 (M2). Điểm ngoặt ≈ 39 \% FSTS (Lind-Mehlum p = .009).

\includegraphics{figures/p3_vietnam/figure_2c.png}

\emph{Hình 2c.} Dự đoán ln(năng suất lao động) như một hàm của FSTS cho
đợt 2023 (M2). Điểm ngoặt ≈ 42 \% FSTS (Lind-Mehlum p = .013).

\includegraphics{figures/p3_vietnam/figure_2d.png}

\emph{Hình 2d.} Dự đoán ln(năng suất lao động) như một hàm của FSTS cho
mẫu gộp (M2). Điểm ngoặt ≈ 40 \% FSTS (Lind-Mehlum p \textless{} .001).

\hypertarget{43-diux1ec5n-giux1ea3i-cuxe1c-kiux1ec3m-ux111ux1ecbnh-giux1ea3-thuyux1ebft}{%
\subsubsection{4.3 Diễn giải các kiểm định giả
thuyết}\label{43-diux1ec5n-giux1ea3i-cuxe1c-kiux1ec3m-ux111ux1ecbnh-giux1ea3-thuyux1ebft}}

H1 được ủng hộ có điều kiện. Kiểm định Lind--Mehlum bác bỏ giả thuyết
không về tính đơn điệu trong cả ba đợt (2009 p = .006, 2015 p = .009,
2023 p = .013) và trong mẫu gộp (p \textless{} .001). Các điểm ngoặt ngụ
ý tập trung chặt chẽ trong khoảng 39,3 \% (2015) đến 46,2 \% (2009) của
cường độ xuất khẩu trực tiếp, với ước lượng gộp tại 39,7 \%. Đồng thời,
các mô hình chỉ riêng doanh nghiệp xuất khẩu (mục 4.5, Panel H) cho thấy
độ cong này yếu đi đáng kể khi biên tham gia được loại trừ. Cách diễn
giải phòng thủ nhất do đó là: chữ U ngược trong toàn mẫu phản ánh cấu
trúc tham gia-cường độ kết hợp, với độ tương phản liên quan đến năng
suất tập trung chủ yếu ở bước chuyển từ không xuất khẩu sang xuất khẩu,
chứ không phải ở độ cong mạnh trong nhóm doanh nghiệp xuất khẩu đơn lẻ.

Ngưỡng gộp 39,7 \% của Việt Nam (phạm vi 39,3--46,2 \% giữa các đợt)
thấp hơn đáng kể so với ngưỡng nhóm phụ sản xuất 47,8 \% được báo cáo
cho doanh nghiệp tư nhân Trung Quốc bởi Author Citation (2026, JFAR).
Khoảng cách này nhất quán với chi phí giao dịch thể chế ràng buộc sớm
hơn trong bối cảnh nền kinh tế mới nổi thu nhập thấp hơn của Việt Nam.
Cụ thể, các doanh nghiệp xuất khẩu Việt Nam tiếp cận chi phí cam kết quá
mức ở mức cường độ xuất khẩu tuyệt đối thấp hơn so với đối tác Trung
Quốc. Tại Trung Quốc, các thể chế thực thi hợp đồng trưởng thành hơn,
thị trường nội địa lớn hơn và hạ tầng hỗ trợ xuất khẩu vững chắc hơn cho
phép doanh nghiệp duy trì cường độ xuất khẩu cao hơn trước khi lợi suất
năng suất giảm dần. Khác biệt ngưỡng giữa các quốc gia do đó cung cấp
bằng chứng giữa các bối cảnh trong WBES cho kênh bối cảnh thể chế trong
các phân tích tổng hợp về mối quan hệ I--P (Wu, Fan, \& Chen, 2022, MIR;
Marano et al., 2016).

H2 được ủng hộ bởi mối liên hệ trực tiếp dương của
\(\tilde{\mathrm{TCI}}\) trong mẫu gộp (β = 0,179, p \textless{} .001)
và trong cả ba đợt (2009 p \textless{} .001; 2015 p = .010; 2023 p =
.006). Mối liên hệ này được củng cố bởi điều tiết TCI có thể phân biệt
thống kê trong ba trong bốn bảng (M3 joint p = .040, .713, .027 và
.003). H3 được ủng hộ tính trung bình nhưng không đồng nhất giữa các
đợt. Ước lượng M7 gộp của \(\tilde{\mathrm{DAI}}\) dương và có ý nghĩa
(β = 0,078, p = .004), nhưng biến đổi sắc nét giữa các đợt: mạnh năm
2009 (β = 0,175, p \textless{} .001), không có ý nghĩa năm 2015 (β =
-0,044, p = .377) và tái xuất hiện năm 2023 (β = 0,095, p = .038). Các
kiểm định z giữa các đợt theo Paternoster (mục 4.5, Panel F) xác nhận
rằng mức giảm 2009 đến 2015 (z = 3,353, p \textless{} .001) và sự phục
hồi 2015 đến 2023 (z = -2,051, p = .040) đều là các dịch chuyển có thể
phân biệt thống kê.

H4 được ủng hộ thăm dò hạn chế. Kiểm định điều tiết liên hợp của DAI
không có ý nghĩa trong năm 2009 (M4 p = .825) và 2015 (M4 p = .125). Vào
năm 2023, kiểm định đạt mức biên ý nghĩa: tương tác riêng lẻ
\(\mathrm{FSTS}^{c}\) × \(\tilde{\mathrm{DAI}}\) = -0,912 (p = .043),
kiểm định liên hợp M4 p = .102, M8 p = .062. Kiểm định liên hợp M8 gộp
(p = .083) cũng ở mức cận ý nghĩa, được chi phối bởi đợt 2023 chứ không
phải bởi điều tiết ổn định xuyên suốt giai đoạn. Kiểm định tương tác
chính thức gộp đợt × biến trọng tâm (Panel I) không phát hiện các khác
biệt giữa các đợt trong các số hạng điều tiết FSTS × DAI. Bằng chứng do
đó được diễn giải như gợi ý về tính có điều kiện đặc trưng theo đợt, chứ
không phải xác nhận về một mẫu hình điều tiết ổn định giữa các đợt. Năm
2023 là đợt duy nhất trong đó điều tiết số có thể phát hiện trong mẫu,
và phát hiện được xem là mang tính thăm dò.

Theo Haans, Pieters, và He (2016), nghiên cứu phân biệt hai loại điều
tiết cho các mối quan hệ I--P cong. Điều tiết loại I làm phẳng hoặc dốc
thêm độ dốc của chữ U ngược (tương tác tuyến tính FSTS × DAI có ý nghĩa
trong khi \(\mathrm{FSTS}^2\) × DAI thì không). Biến điều tiết loại I
dịch chuyển vị trí của điểm ngoặt nhưng giữ lại hình chữ U ngược. Điều
tiết loại II lật hình dạng của đường cong (\(\mathrm{FSTS}^2\) × DAI có
ý nghĩa): giá trị biến điều tiết cao so với thấp tạo ra các dạng hàm
khác biệt về chất. Bằng chứng DAI năm 2023 phù hợp với Loại I. Cụ thể,
tương tác tuyến tính có ý nghĩa \(\mathrm{FSTS}^{c}\) ×
\(\tilde{\mathrm{DAI}}\) = −0,912 (p = .043) cho thấy áp dụng công nghệ
số làm suy yếu độ dốc dương ở cường độ xuất khẩu thấp, kéo điểm ngoặt
vào trong. Trong khi đó, tương tác \(\mathrm{FSTS}^2\) × DAI vẫn không
có ý nghĩa, xác nhận rằng hình chữ U ngược được giữ lại. Điều tiết Loại
II (năng lực số đảo ngược độ cong cho các doanh nghiệp áp dụng cao)
không được ủng hộ trong mẫu Việt Nam.

\includegraphics{figures/p3_vietnam/figure_3_moderator_marginals.png}

\emph{Hình 3.} Tác động biên của \(\tilde{\mathrm{TCI}}\) và
\(\tilde{\mathrm{DAI}}\) lên ln(năng suất lao động) qua các mức FSTS
(M7/M8). TCI cho thấy một dịch chuyển mức dương ổn định giữa các đợt;
điều tiết DAI đặc trưng theo đợt, với tương tác năm 2023 cho thấy suy
giảm ở FSTS cao (Điều tiết loại I: làm phẳng độ dốc, giữ nguyên hình
dạng).

\hypertarget{44-mux1eabu-huxecnh-thux1ef1c-nghiux1ec7m-chuxednh-tham-gia--cux1b0ux1eddng-ux111ux1ed9}{%
\subsubsection{4.4 Mẫu hình thực nghiệm chính: tham gia × cường
độ}\label{44-mux1eabu-huxecnh-thux1ef1c-nghiux1ec7m-chuxednh-tham-gia--cux1b0ux1eddng-ux111ux1ed9}}

Trước khi đọc bảng tiếp theo, người đọc cần được neo trong cấu trúc hai
biên ở mục 2.1. Chữ U ngược trong toàn mẫu cung cấp thông tin về mẫu
hình tham gia-cường độ kết hợp, nhưng độ cong của nó được nhận diện chủ
yếu thông qua biên tham gia. Cụ thể, chỉ khoảng 1,0 \% doanh nghiệp gộp
nằm trong ±5 điểm phần trăm của các điểm ngoặt đặc trưng theo đợt (xem
kiểm tra mật độ trong mục 4.5), và phần lớn khối lượng nằm ở FSTS = 0.
Khi M2/M7/M8 được tái ước lượng trên nhóm phụ chỉ doanh nghiệp xuất khẩu
(FSTS \textgreater{} 0; gộp N = 669, mục 4.5 Panel H), số hạng tuyến
tính \(\mathrm{FSTS}^{c}\) âm (β = -0,861, p \textless{} .001) nhưng số
hạng bậc hai không có ý nghĩa (\((\mathrm{FSTS}^{c})^2\) β = -0,200, p =
.660, M8 joint p = .462). H1a (biên tham gia) do đó là biên liên quan
đến năng suất chủ đạo trong bộ dữ liệu này. Độ cong cường độ trong nhóm
doanh nghiệp xuất khẩu mà H1b tuyên bố yếu đi khi biên tham gia được
loại trừ. Chữ U ngược trong toàn mẫu được giữ lại như tính đều đặn thực
nghiệm tiêu điểm, nhưng được diễn giải rõ ràng qua lăng kính cơ chế kép:
phần lớn chênh lệch năng suất nằm giữa doanh nghiệp không xuất khẩu và
doanh nghiệp xuất khẩu, với độ cong bổ sung hạn chế trong nhóm doanh
nghiệp xuất khẩu. Bảng 1 do đó đọc như mô tả về mẫu hình tham gia ×
cường độ kết hợp, chứ không phải tuyên bố cấu trúc về độ cong cường độ
trong nhóm doanh nghiệp xuất khẩu.

Bảng 1 tóm tắt diễn giải định hướng của các hệ số trọng tâm từ các đặc
tả đầy đủ theo từng đợt và cho mẫu gộp.

\hypertarget{45-kiux1ec3m-tra-ux111ux1ed9-vux1eefng}{%
\subsubsection{4.5 Kiểm tra độ
vững}\label{45-kiux1ec3m-tra-ux111ux1ed9-vux1eefng}}

Bốn nhóm kiểm tra độ vững xem xét liệu các suy luận trung tâm có nhạy
cảm với các lựa chọn đo lường, thành phần ngành, cấu trúc tham gia xuất
khẩu và nội sinh hay không. Các bảng bên dưới được ước lượng trên cùng
thiết kế OLS HC1 như các mô hình chính.

\textbf{Panel G --- Tách ngành (sản xuất so với phi sản xuất).} Để kiểm
tra liệu các kênh số và năng lực có vận hành đồng nhất qua thành phần
ngành rộng của nhóm doanh nghiệp xuất khẩu Việt Nam hay không, nghiên
cứu tái ước lượng các đặc tả M2/M7/M8 gộp riêng cho doanh nghiệp sản
xuất (sector1 ∈ \{1, 2, 3\}, ISIC 15--37, N = 1.854) và doanh nghiệp phi
sản xuất (tiện ích, xây dựng, bán buôn và bán lẻ, vận tải, tài chính và
các dịch vụ khác; sector1 ∈ \{4, 5, 6, 7\}; N = 1.104).

Chữ U ngược được giữ lại trong cả hai nhóm phụ nhưng sắc nét hơn ở sản
xuất (\(\mathrm{FSTS}^{c}\) β = 0,971, p = .001;
\((\mathrm{FSTS}^{c})^2\) β = -1,883, p \textless{} .001) so với phi sản
xuất (\(\mathrm{FSTS}^{c}\) β = 1,615, p = .064, cận ngưỡng;
\((\mathrm{FSTS}^{c})^2\) β = -2,479, p = .046).

Quan trọng hơn về mặt thực chất, các kênh năng lực và áp dụng công nghệ
số vận hành chủ yếu ở sản xuất.

Trong đặc tả M7 trực tiếp kép, doanh nghiệp sản xuất cho thấy mối liên
hệ trực tiếp \(\tilde{\mathrm{TCI}}\) mạnh (β = 0,223, p \textless{}
.001) và mối liên hệ trực tiếp \(\tilde{\mathrm{DAI}}\) dương (β =
0,087, p = .009). Trong khi đó, doanh nghiệp phi sản xuất chỉ cho thấy
tác động \(\tilde{\mathrm{TCI}}\) cận ngưỡng (β = 0,090, p = .096) và
tác động \(\tilde{\mathrm{DAI}}\) không có ý nghĩa (β = 0,068, p =
.133). Kênh điều tiết DAI tập trung ở sản xuất (M8 joint p = .103, cận
ngưỡng; \(\mathrm{FSTS}^{c}\) × \(\tilde{\mathrm{DAI}}\) = -0,543, p =
.079) và không có ý nghĩa đồng đều ở phi sản xuất (M8 joint p = .280).
Ngược lại, điều tiết TCI có thể phân biệt thống kê với không trong cả
hai nhóm phụ (M3 joint p = .011 ở sản xuất và p = .007 ở phi sản xuất).
Kết quả cho thấy tác động làm phẳng độ cong của năng lực hoạt động rộng
rãi, trong khi kênh số vẫn đặc thù cho cơ sở doanh nghiệp xuất khẩu sản
xuất. Về mặt thực chất, mẫu hình này nhất quán với quan điểm rằng áp
dụng công nghệ số cơ bản tương tác với các yêu cầu phối hợp sản xuất
xuyên biên giới mạnh hơn ở các ngành hàng hoá có thể giao dịch, so với
các ngành định hướng dịch vụ hoặc định hướng nội địa của nền kinh tế
Việt Nam.

\textbf{Panel H --- Nhóm phụ chỉ doanh nghiệp xuất khẩu (FSTS
\textgreater{} 0).} Tỉ lệ doanh nghiệp xuất khẩu là 28,4 \%, 20,7 \% và
18,8 \% qua các đợt. Vì vậy, chữ U ngược khớp trên toàn mẫu được nhận
diện một phần bởi biên tham gia giữa FSTS = 0 và FSTS \textgreater{} 0.
Nghiên cứu tái khớp M2/M7/M8 trên nhóm phụ chỉ doanh nghiệp xuất khẩu (N
≈ 281 năm 2009, 198 năm 2015, 190 năm 2023, và 669 gộp). Đặc tả chỉ
riêng doanh nghiệp xuất khẩu gộp cho số hạng tuyến tính
\(\mathrm{FSTS}^{c}\) âm (β = -0,861, p \textless{} .001) nhưng số hạng
bậc hai không có ý nghĩa (\((\mathrm{FSTS}^{c})^2\) β = -0,200, p =
.660). Kiểm định M8 liên hợp của khối độ cong cộng điều tiết không có ý
nghĩa (joint F p = .462). Các ước lượng chỉ riêng doanh nghiệp xuất khẩu
theo từng đợt cũng nhiễu hơn và yếu hơn riêng lẻ so với các đối tác toàn
mẫu.

Nghiên cứu đọc kết quả này như chứng minh rằng chữ U ngược trong đặc tả
chính được nhận diện có ý nghĩa bởi biên tham gia, chứ không hoàn toàn
bởi biến thiên cường độ trong nhóm doanh nghiệp xuất khẩu. Phát hiện
nhất quán với bằng chứng mô tả ở mục 3.1: cường độ xuất khẩu trực tiếp
là biến lạm phát ở mức không trong dữ liệu WBES Việt Nam, với khối lượng
hạn chế giữa biên tham gia và điểm ngoặt ngụ ý. Tuyên bố H1 thực chất do
đó được đọc tốt nhất như mối liên hệ không đơn điệu giữa tham
gia-và-cường độ trong xuất khẩu và năng suất, chứ không phải như tuyên
bố chặt chẽ về độ cong cường độ trong nhóm doanh nghiệp xuất khẩu.

\textbf{Panel I --- Kiểm định tương tác đợt × biến trọng tâm gộp.} Để
kiểm định chính thức liệu các mẫu hình độ cong và điều tiết theo từng
đợt có thể tách biệt thống kê với các ước lượng gộp hay không, nghiên
cứu tái ước lượng M8 gộp với tập tương tác đợt bão hoà trên các số hạng
trọng tâm (\(\mathrm{FSTS}^{c}\) × đợt, \((\mathrm{FSTS}^{c})^2\) × đợt,
\(\tilde{\mathrm{DAI}}\) × đợt, \(\tilde{\mathrm{TCI}}\) × đợt). Các
kiểm định Wald liên hợp cho thấy chỉ \(\tilde{\mathrm{DAI}}\) × đợt là
có thể phân biệt thống kê với mức trung bình gộp (joint p = .016). Các
khác biệt tác động trực tiếp giữa các đợt của \(\mathrm{FSTS}^{c}\),
\((\mathrm{FSTS}^{c})^2\) và \(\tilde{\mathrm{TCI}}\) không thể tách
biệt thống kê (tất cả joint p \textgreater{} .25). Các khác biệt giữa
các đợt của \(\mathrm{FSTS}^{c}\) × \(\tilde{\mathrm{DAI}}\) và
\((\mathrm{FSTS}^{c})^2\) × \(\tilde{\mathrm{DAI}}\) cũng không thể tách
biệt (joint p \textgreater{} .55). Kiểm định chính thức này xác nhận kết
quả Paternoster mô tả: chỉ các dịch chuyển trực tiếp DAI là có thể phân
biệt giữa các đợt, trong khi các tham số độ cong và các số hạng điều
tiết FSTS × DAI chia sẻ độ lớn gộp chung giữa các đợt.

\textbf{Kiểm tra mật độ quanh điểm ngoặt.} Chúng tôi cũng báo cáo khối
lượng trong mẫu của doanh nghiệp trong vùng độ cong. Xác định một dải ±5
điểm phần trăm quanh điểm ngoặt đặc trưng theo đợt (TP), tỉ lệ doanh
nghiệp với FSTS trong dải là 0,6 \% (6 trong 989) năm 2009, 1,2 \% (11
trong 956) năm 2015, 0,9 \% (9 trong 1.013) năm 2023 và 1,0 \% (29 trong
2.958) gộp. Phần lớn khối lượng nằm ở FSTS = 0; các ước lượng điểm ngoặt
được nhận diện chủ yếu thông qua độ tương phản giữa doanh nghiệp không
xuất khẩu và doanh nghiệp xuất khẩu và đuôi phải phía trên TP, chứ không
phải thông qua biến thiên dày đặc trong dải.

Người đọc nên ghi nhớ cấu trúc mật độ này khi diễn giải độ lớn chính xác
của điểm ngoặt.

\hypertarget{ux111ux1ed9-vux1eefng-ux111ux1ed1i-vux1edbi-nux1ed9i-sinh-vuxe0-chux1ecdn-lux1ef1a}{%
\paragraph{Độ vững đối với nội sinh và chọn
lựa}\label{ux111ux1ed9-vux1eefng-ux111ux1ed1i-vux1edbi-nux1ed9i-sinh-vuxe0-chux1ecdn-lux1ef1a}}

Ba cách tiếp cận bổ sung giải quyết các quan ngại về chọn lựa và nội
sinh. Hiệu chỉnh Heckman hai bước (Panel E) cho các tỉ số Mills nghịch
đảo (inverse Mills ratio, IMR) không có ý nghĩa qua tất cả các đợt (tất
cả \textbar λ\textbar{} \textless{} 0,84, p \textgreater{} .25), cho
thấy không có thiên lệch chọn lựa có thể phát hiện trên nhóm phụ doanh
nghiệp xuất khẩu. Đối sánh điểm xu hướng (propensity-score matching,
PSM; Panel J) xác nhận các mối liên hệ trung bình H2 và H3 mà không áp
đặt tính tuyến tính: các ước lượng ATT cho sở hữu website (0,298--0,321,
p \textless{} .001) và cho trạng thái công nghệ nước ngoài / chứng nhận
(0,609--0,637, p \textless{} .001) nhất quán về mặt thực chất với các
kết quả OLS. Các chẩn đoán cân bằng đối sánh đầy đủ và các ước lượng ATT
được báo cáo trong Bảng A của Phụ lục Trực tuyến.

Ước lượng bình phương nhỏ nhất hai bước (two-stage least squares, 2SLS)
sử dụng tỉ lệ áp dụng đồng cấp ngành × vùng × đợt rút một (Panel K;
thống kê F giai đoạn một 22--35, vượt xa ngưỡng Staiger--Stock) cho ước
lượng DAI được công cụ hoá là 0,018 (p = .942) và ước lượng TCI được
công cụ hoá là 1,639 (p \textless{} .001). Hệ số DAI 2SLS không có ý
nghĩa là một kết quả độ vững then chốt: nó xác nhận rằng hiện diện
website không gây ra cải thiện năng suất một cách khả tín ở Việt Nam
2023, củng cố cách diễn giải tính lỗi thời của biến đại diện Tầng 1 cho
tương tác DAI×FSTS âm. Hệ số TCI được công cụ hoá dương mạnh (β=1,639,
p\textless.001, lớn hơn đáng kể so với ước lượng OLS, nhất quán với hiệu
chỉnh thiên lệch suy giảm) củng cố thêm rằng năng lực công nghệ nước
ngoài và tiêu chuẩn là cơ chế liên quan đến năng suất vững chắc trong
bối cảnh thể chế này. Tập công cụ (tỉ lệ áp dụng đồng cấp trong các ô
ngành × vùng × đợt) thoả mãn tính liên quan (giai đoạn một mạnh) và loại
trừ xấp xỉ (tỉ lệ áp dụng đồng cấp khó có khả năng ảnh hưởng đến năng
suất doanh nghiệp riêng lẻ qua các kênh khác ngoài áp dụng DAI, có điều
kiện trên các ô ngành-đợt). Các giới hạn ổn định δ của Oster (2019) (giả
định R\_max = 1,3 × R\_controlled) xác nhận rằng không có hệ số trọng
tâm nào thay đổi dấu hoặc sụp đổ về không dưới các độ lớn khả tín của
chọn lựa không quan sát. Kết quả giai đoạn một 2SLS đầy đủ và các tính
toán giới hạn Oster được báo cáo trong Bảng B của Phụ lục Trực tuyến.

Các kiểm tra nhạy cảm đo lường xác nhận rằng các suy luận cốt lõi không
phụ thuộc vào việc xây dựng chỉ số hợp thành. Làm giàu TCI với các chỉ
báo R\&D và đổi mới sản phẩm (Panel A) làm suy giảm hệ số TCI khiêm tốn
nhưng không thay đổi các kết luận về chất. Chỉ số hợp thành
\(\mathrm{DAI}_{\text{rich}}\) được làm giàu (Tầng 1--2, Panel B) cho
điều tiết theo cùng hướng vào năm 2023 với ý nghĩa yếu hơn. Loại trừ
doanh nghiệp siêu nhỏ (Panel D) bảo tồn chữ U ngược và các mối liên hệ
TCI. Các kiểm định z giữa các đợt theo Paternoster et al. (Panel F) xác
nhận rằng dịch chuyển vòng đời DAI có thể phân biệt thống kê trong khi
các tham số độ cong chia sẻ một độ lớn gộp chung. Các ước lượng cấp bảng
được báo cáo trong các Bảng C--F của Phụ lục Trực tuyến.

\textbf{Lưu ý về kiểm định bội.} Mục 4.5 báo cáo bốn bảng tường thuật
(G, H, I, F) và ba nhóm kiểm tra độ vững (nội sinh/chọn lựa, nhạy cảm đo
lường) qua nhiều số hạng trọng tâm. Chúng tôi không áp dụng hiệu chỉnh
kiểm định bội chính thức vì các bảng kiểm tra các quan ngại nhận diện
khác nhau thay vì kiểm định cùng một giả thuyết lặp đi lặp lại, nhưng
người đọc nên cân nhắc bất kỳ kết quả bảng cận ngưỡng đơn lẻ nào tương
ứng. Các suy luận thực chất của chúng tôi trong Mục 5 dựa vào mẫu hình
qua các bảng và tính nhất quán định hướng của các ước lượng trọng tâm
chứ không dựa vào ý nghĩa của bất kỳ bảng kiểm tra độ vững đơn lẻ nào.

Bảng 3 tổng hợp các bảng kiểm tra độ vững được ghi nhận trong mục 4.5
trong một tổng quan duy nhất để dễ so sánh chéo.

Bảng 4 báo cáo các điểm ngoặt ngụ ý của đặc tả chữ U ngược (M2) và các
giá trị p Lind--Mehlum cho mỗi đợt và mẫu gộp.

Bảng 3: Các bảng kiểm tra độ vững A--K và các kiểm tra bổ sung.

\begin{longtable}[]{@{}llllll@{}}
\toprule\noalign{}
Bảng & Mẫu & N & Số hạng & b (SE) & p \\
\midrule\noalign{}
\endhead
\bottomrule\noalign{}
\endlastfoot
DAI rich DAI rich bin & 2023 & 1013 & FSTSc & +1,002† (0,607) & 0,099 \\
DAI rich DAI rich bin & 2023 & 1013 & FSTSc2 & -1,707* (0,813) &
0,036 \\
DAI rich DAI rich bin & 2023 & 1013 & \(\tilde{\mathrm{TCI}}\) &
+0,129** (0,045) & 0,004 \\
DAI rich DAI rich bin & 2023 & 1013 & DAI\_rich\_bin\_z & -0,009 (0,100)
& 0,926 \\
DAI rich DAI rich bin & 2023 & 1013 & FSTSc\_DAI\_rich\_bin & -0,893
(0,649) & 0,169 \\
DAI rich DAI rich bin & 2023 & 1013 & FSTSc2\_DAI\_rich\_bin & +1,038
(0,878) & 0,237 \\
DAI rich DAI rich bin & 2023 & 1013 & joint\_F\_DAI\_rich\_interactions
& +1,4181 & 0,243 \\
DAI rich DAI rich cont & 2023 & 1013 & FSTSc & +1,032† (0,569) &
0,070 \\
DAI rich DAI rich cont & 2023 & 1013 & FSTSc2 & -1,744* (0,757) &
0,021 \\
DAI rich DAI rich cont & 2023 & 1013 & \(\tilde{\mathrm{TCI}}\) &
+0,126** (0,045) & 0,005 \\
DAI rich DAI rich cont & 2023 & 1013 & DAI\_rich\_cont\_z & -0,004
(0,083) & 0,957 \\
DAI rich DAI rich cont & 2023 & 1013 & FSTSc\_DAI\_rich\_cont & -0,933†
(0,526) & 0,076 \\
DAI rich DAI rich cont & 2023 & 1013 & FSTSc2\_DAI\_rich\_cont & +1,052
(0,722) & 0,145 \\
DAI rich DAI rich cont & 2023 & 1013 & joint\_F\_DAI\_rich\_interactions
& +2,3148† & 0,099 \\
DAI thin on rich sample 2023 & 2023 & 1013 & FSTSc & +1,072* (0,493) &
0,030 \\
DAI thin on rich sample 2023 & 2023 & 1013 & FSTSc2 & -1,793** (0,666) &
0,007 \\
DAI thin on rich sample 2023 & 2023 & 1013 & \(\tilde{\mathrm{TCI}}\) &
+0,129** (0,045) & 0,004 \\
DAI thin on rich sample 2023 & 2023 & 1013 & \(\tilde{\mathrm{DAI}}\) &
-0,011 (0,073) & 0,880 \\
DAI thin on rich sample 2023 & 2023 & 1013 & FSTSc\_DAIz & -0,912*
(0,450) & 0,043 \\
DAI thin on rich sample 2023 & 2023 & 1013 & FSTSc2\_DAIz & +1,043†
(0,633) & 0,100 \\
DAI thin on rich sample 2023 & 2023 & 1013 &
joint\_F\_DAI\_thin\_interactions & +2,7833† & 0,062 \\
IV 2SLS K DAI & gộp & 2298 & \(\tilde{\mathrm{DAI}}\) (công cụ hoá) &
+0,018 (0,249) & 0,942 \\
IV 2SLS K TCI & gộp & 2298 & \(\tilde{\mathrm{TCI}}\) (công cụ hoá) &
+1,639*** (0,299) & 0,000 \\
PSM J NN1 caliper005 & gộp & 1085 & ATT\_DAI\_treat & +0,298*** (0,061)
& 0,000 \\
PSM J NN1 caliper005 & gộp & 644 & ATT\_TCI\_treat & +0,637*** (0,077) &
0,000 \\
PSM J kernel bw006 & gộp & 1085 & ATT\_DAI\_treat & +0,321*** (0,043) &
0,000 \\
PSM J kernel bw006 & gộp & 639 & ATT\_TCI\_treat & +0,609*** (0,056) &
0,000 \\
TCI full direct & 2015 & 956 & FSTSc & +1,139* (0,539) & 0,035 \\
TCI full direct & 2015 & 956 & FSTSc2 & -2,088** (0,748) & 0,005 \\
TCI full direct & 2015 & 956 & TCI\_full\_z & +0,056 (0,049) & 0,253 \\
TCI full direct & 2015 & 956 & \(\tilde{\mathrm{DAI}}\) & -0,033 (0,051)
& 0,520 \\
TCI full direct & 2023 & 1013 & FSTSc & +0,675 (0,475) & 0,155 \\
TCI full direct & 2023 & 1013 & FSTSc2 & -1,266* (0,646) & 0,050 \\
TCI full direct & 2023 & 1013 & TCI\_full\_z & +0,096* (0,043) &
0,024 \\
TCI full direct & 2023 & 1013 & \(\tilde{\mathrm{DAI}}\) & +0,097*
(0,046) & 0,037 \\
TCI full moderation & 2015 & 956 & FSTSc & +1,347* (0,598) & 0,024 \\
TCI full moderation & 2015 & 956 & FSTSc2 & -2,350** (0,829) & 0,005 \\
TCI full moderation & 2015 & 956 & TCI\_full\_z & +0,004 (0,072) &
0,960 \\
TCI full moderation & 2015 & 956 & FSTSc\_TCIfull & -0,545 (0,495) &
0,270 \\
TCI full moderation & 2015 & 956 & FSTSc2\_TCIfull & +0,641 (0,694) &
0,356 \\
TCI full moderation & 2015 & 956 & joint\_F\_TCI\_full\_interactions &
+0,7096 & 0,492 \\
TCI full moderation & 2023 & 1013 & FSTSc & +0,949† (0,574) & 0,098 \\
TCI full moderation & 2023 & 1013 & FSTSc2 & -1,579* (0,777) & 0,042 \\
TCI full moderation & 2023 & 1013 & TCI\_full\_z & +0,117* (0,049) &
0,016 \\
TCI full moderation & 2023 & 1013 & FSTSc\_TCIfull & -0,317 (0,314) &
0,314 \\
TCI full moderation & 2023 & 1013 & FSTSc2\_TCIfull & +0,185 (0,431) &
0,668 \\
TCI full moderation & 2023 & 1013 & joint\_F\_TCI\_full\_interactions &
+1,8868 & 0,152 \\
common N reconciled 2023 & M8 & 1013 & FSTSc & +1,072* (0,493) &
0,030 \\
common N reconciled 2023 & M8 & 1013 & FSTSc2 & -1,793** (0,666) &
0,007 \\
common N reconciled 2023 & M8 & 1013 & \(\tilde{\mathrm{TCI}}\) &
+0,129** (0,045) & 0,004 \\
common N reconciled 2023 & M8 & 1013 & \(\tilde{\mathrm{DAI}}\) & -0,011
(0,073) & 0,880 \\
common N reconciled 2023 & M8 & 1013 & FSTSc\_DAIz & -0,912* (0,450) &
0,043 \\
common N reconciled 2023 & M8 & 1013 & FSTSc2\_DAIz & +1,043† (0,633) &
0,100 \\
common N reconciled 2023 & 2023 & 1013 & joint\_F\_DAI\_M8 & +2,7833† &
0,062 \\
exporter only & M2 & 281 & FSTSc & -0,611* (0,242) & 0,012 \\
exporter only & M2 & 281 & FSTSc2 & +1,674* (0,853) & 0,050 \\
exporter only & M7 & 281 & FSTSc & -0,409† (0,238) & 0,085 \\
exporter only & M7 & 281 & FSTSc2 & +1,585† (0,854) & 0,064 \\
exporter only & M7 & 281 & \(\tilde{\mathrm{TCI}}\) & +0,149† (0,077) &
0,053 \\
exporter only & M7 & 281 & \(\tilde{\mathrm{DAI}}\) & +0,192* (0,090) &
0,034 \\
exporter only & M8 & 281 & FSTSc & -0,403 (0,262) & 0,124 \\
exporter only & M8 & 281 & FSTSc2 & +1,377 (0,918) & 0,134 \\
exporter only & M8 & 281 & \(\tilde{\mathrm{TCI}}\) & +0,147† (0,077) &
0,057 \\
exporter only & M8 & 281 & \(\tilde{\mathrm{DAI}}\) & +0,075 (0,143) &
0,598 \\
exporter only & M8 & 281 & FSTSc\_DAIz & +0,078 (0,238) & 0,743 \\
exporter only & M8 & 281 & FSTSc2\_DAIz & +0,802 (0,874) & 0,359 \\
exporter only & 2009\_exp & 281 & joint\_F\_DAI\_M8 & +0,4235 & 0,655 \\
exporter only & M2 & 198 & FSTSc & -0,773** (0,290) & 0,008 \\
exporter only & M2 & 198 & FSTSc2 & -2,204* (0,957) & 0,021 \\
exporter only & M7 & 198 & FSTSc & -0,771** (0,296) & 0,009 \\
exporter only & M7 & 198 & FSTSc2 & -2,619** (0,957) & 0,006 \\
exporter only & M7 & 198 & \(\tilde{\mathrm{TCI}}\) & +0,242*** (0,071)
& 0,001 \\
exporter only & M7 & 198 & \(\tilde{\mathrm{DAI}}\) & -0,212† (0,112) &
0,059 \\
exporter only & M8 & 198 & FSTSc & -0,672* (0,301) & 0,025 \\
exporter only & M8 & 198 & FSTSc2 & -2,719* (1,147) & 0,018 \\
exporter only & M8 & 198 & \(\tilde{\mathrm{TCI}}\) & +0,237*** (0,071)
& 0,001 \\
exporter only & M8 & 198 & \(\tilde{\mathrm{DAI}}\) & -0,177 (0,202) &
0,383 \\
exporter only & M8 & 198 & FSTSc\_DAIz & -0,335 (0,277) & 0,226 \\
exporter only & M8 & 198 & FSTSc2\_DAIz & -0,172 (1,096) & 0,876 \\
exporter only & 2015\_exp & 198 & joint\_F\_DAI\_M8 & +0,7344 & 0,481 \\
exporter only & M2 & 190 & FSTSc & -0,910** (0,335) & 0,007 \\
exporter only & M2 & 190 & FSTSc2 & -1,461 (1,160) & 0,208 \\
exporter only & M7 & 190 & FSTSc & -0,856* (0,344) & 0,013 \\
exporter only & M7 & 190 & FSTSc2 & -1,397 (1,171) & 0,233 \\
exporter only & M7 & 190 & \(\tilde{\mathrm{TCI}}\) & +0,061 (0,067) &
0,366 \\
exporter only & M7 & 190 & \(\tilde{\mathrm{DAI}}\) & +0,037 (0,086) &
0,665 \\
exporter only & M8 & 190 & FSTSc & -1,079** (0,332) & 0,001 \\
exporter only & M8 & 190 & FSTSc2 & -2,921** (1,094) & 0,008 \\
exporter only & M8 & 190 & \(\tilde{\mathrm{TCI}}\) & +0,080 (0,067) &
0,234 \\
exporter only & M8 & 190 & \(\tilde{\mathrm{DAI}}\) & -0,271† (0,141) &
0,054 \\
exporter only & M8 & 190 & FSTSc\_DAIz & +0,510† (0,294) & 0,083 \\
exporter only & M8 & 190 & FSTSc2\_DAIz & +2,676** (1,021) & 0,009 \\
exporter only & 2023\_exp & 190 & joint\_F\_DAI\_M8 & +3,4455* &
0,034 \\
exporter only & M2 & 669 & FSTSc & -0,861*** (0,167) & 0,000 \\
exporter only & M2 & 669 & FSTSc2 & -0,200 (0,581) & 0,730 \\
exporter only & M7 & 669 & FSTSc & -0,704*** (0,173) & 0,000 \\
exporter only & M7 & 669 & FSTSc2 & -0,172 (0,577) & 0,765 \\
exporter only & M7 & 669 & \(\tilde{\mathrm{TCI}}\) & +0,178*** (0,042)
& 0,000 \\
exporter only & M7 & 669 & \(\tilde{\mathrm{DAI}}\) & +0,066 (0,057) &
0,247 \\
exporter only & M8 & 669 & FSTSc & -0,683*** (0,183) & 0,000 \\
exporter only & M8 & 669 & FSTSc2 & -0,497 (0,643) & 0,440 \\
exporter only & M8 & 669 & \(\tilde{\mathrm{TCI}}\) & +0,179*** (0,042)
& 0,000 \\
exporter only & M8 & 669 & \(\tilde{\mathrm{DAI}}\) & -0,025 (0,104) &
0,809 \\
exporter only & M8 & 669 & FSTSc\_DAIz & -0,001 (0,164) & 0,995 \\
exporter only & M8 & 669 & FSTSc2\_DAIz & +0,702 (0,625) & 0,261 \\
exporter only & pooled\_exp & 669 & joint\_F\_DAI\_M8 & +0,7733 &
0,462 \\
micro excluded pooled & pooled\_l1ge10 & 2473 & FSTSc & +0,794** (0,301)
& 0,008 \\
micro excluded pooled & pooled\_l1ge10 & 2473 & FSTSc2 & -1,647***
(0,434) & 0,000 \\
micro excluded pooled & pooled\_l1ge10 & 2473 & \(\tilde{\mathrm{TCI}}\)
& +0,188*** (0,028) & 0,000 \\
micro excluded pooled & pooled\_l1ge10 & 2473 & \(\tilde{\mathrm{DAI}}\)
& +0,054 (0,051) & 0,288 \\
micro excluded pooled & pooled\_l1ge10 & 2473 & FSTSc\_DAIz & -0,387
(0,281) & 0,169 \\
micro excluded pooled & pooled\_l1ge10 & 2473 & FSTSc2\_DAIz & +0,405
(0,416) & 0,330 \\
micro excluded pooled & pooled\_l1ge10 & 2473 &
joint\_F\_DAI\_interactions & +1,7902 & 0,167 \\
sector split manufacturing & M2 & 1854 & FSTSc & +0,971** (0,296) &
0,001 \\
sector split manufacturing & M2 & 1854 & FSTSc2 & -1,883*** (0,426) &
0,000 \\
sector split manufacturing & M7 & 1854 & FSTSc & +0,662* (0,296) &
0,025 \\
sector split manufacturing & M7 & 1854 & FSTSc2 & -1,382** (0,428) &
0,001 \\
sector split manufacturing & M7 & 1854 & \(\tilde{\mathrm{TCI}}\) &
+0,223*** (0,033) & 0,000 \\
sector split manufacturing & M7 & 1854 & \(\tilde{\mathrm{DAI}}\) &
+0,087** (0,033) & 0,009 \\
sector split manufacturing & M8 & 1854 & FSTSc & +0,840* (0,328) &
0,010 \\
sector split manufacturing & M8 & 1854 & FSTSc2 & -1,624*** (0,473) &
0,001 \\
sector split manufacturing & M8 & 1854 & \(\tilde{\mathrm{TCI}}\) &
+0,228*** (0,033) & 0,000 \\
sector split manufacturing & M8 & 1854 & \(\tilde{\mathrm{DAI}}\) &
+0,034 (0,057) & 0,552 \\
sector split manufacturing & M8 & 1854 & FSTSc\_DAIz & -0,543† (0,310) &
0,079 \\
sector split manufacturing & M8 & 1854 & FSTSc2\_DAIz & +0,626 (0,457) &
0,170 \\
sector split manufacturing & gộp & 1854 &
joint\_F\_DAI\_interactions\_M8 & +2,2806 & 0,102 \\
sector split manufacturing & gộp & 1854 &
joint\_F\_TCI\_interactions\_M3 & +4,5710* & 0,011 \\
sector split non manufacturing & M2 & 1104 & FSTSc & +1,615† (0,870) &
0,064 \\
sector split non manufacturing & M2 & 1104 & FSTSc2 & -2,479* (1,244) &
0,046 \\
sector split non manufacturing & M7 & 1104 & FSTSc & +1,350 (0,890) &
0,130 \\
sector split non manufacturing & M7 & 1104 & FSTSc2 & -2,039 (1,275) &
0,110 \\
sector split non manufacturing & M7 & 1104 & \(\tilde{\mathrm{TCI}}\) &
+0,090† (0,054) & 0,096 \\
sector split non manufacturing & M7 & 1104 & \(\tilde{\mathrm{DAI}}\) &
+0,068 (0,045) & 0,133 \\
sector split non manufacturing & M8 & 1104 & FSTSc & +1,387 (0,868) &
0,110 \\
sector split non manufacturing & M8 & 1104 & FSTSc2 & -2,115† (1,237) &
0,087 \\
sector split non manufacturing & M8 & 1104 & \(\tilde{\mathrm{TCI}}\) &
+0,091† (0,054) & 0,091 \\
sector split non manufacturing & M8 & 1104 & \(\tilde{\mathrm{DAI}}\) &
+0,060 (0,122) & 0,622 \\
sector split non manufacturing & M8 & 1104 & FSTSc\_DAIz & -0,190
(0,719) & 0,791 \\
sector split non manufacturing & M8 & 1104 & FSTSc2\_DAIz & -0,341
(1,065) & 0,749 \\
sector split non manufacturing & gộp & 1104 &
joint\_F\_DAI\_interactions\_M8 & +1,2729 & 0,280 \\
sector split non manufacturing & gộp & 1104 &
joint\_F\_TCI\_interactions\_M3 & +4,9441** & 0,007 \\
wave interaction pooled & gộp & 2958 & joint\_F\_FSTSc\_x\_wave &
+0,8901 & 0,411 \\
wave interaction pooled & gộp & 2958 & joint\_F\_FSTSc2\_x\_wave &
+0,7756 & 0,461 \\
wave interaction pooled & gộp & 2958 & joint\_F\_DAIz\_x\_wave &
+4,1656* & 0,016 \\
wave interaction pooled & gộp & 2958 & joint\_F\_TCIz\_x\_wave & +1,1156
& 0,328 \\
wave interaction pooled & gộp & 2958 & joint\_F\_all\_wave\_interactions
& +2,0664* & 0,036 \\
\end{longtable}

\emph{Ghi chú.} Các bảng kiểm tra độ vững:
\(\mathrm{TCI}_{\text{full}}\) = chỉ số hợp thành TCI rộng hơn;
\(\mathrm{DAI}_{\text{rich}}\) = chỉ số hợp thành DAI giàu hơn (chỉ
2023). Sai số chuẩn vững HC1. † p \textless{} .10; * p \textless{} .05;
** p \textless{} .01; *** p \textless{} .001.

Bảng 4: Các điểm ngoặt ngụ ý của chữ U ngược (đặc tả M2).

\begin{longtable}[]{@{}lllll@{}}
\toprule\noalign{}
Mẫu & Điểm ngoặt (FSTS) & 95\% CI & Lind-Mehlum p & Phạm vi FSTS \\
\midrule\noalign{}
\endhead
\bottomrule\noalign{}
\endlastfoot
2009 (N=989) & 46,2\% & {[}37,4\%, 55,1\%{]} & 0,006** & {[}0\%,
100\%{]} \\
2015 (N=956) & 39,3\% & {[}30,3\%, 48,4\%{]} & 0,009** & {[}0\%,
100\%{]} \\
2023 (N=1.013) & 41,6\% & {[}31,7\%, 51,5\%{]} & 0,013* & {[}0\%,
100\%{]} \\
Gộp (N=2.958) & 39,7\% & {[}34,0\%, 45,5\%{]} & 0,000*** & {[}0\%,
100\%{]} \\
\end{longtable}

\emph{Ghi chú.} Các điểm ngoặt được ước lượng từ M2 (đặc tả FSTS bậc
hai). Giá trị p của kiểm định Lind-Mehlum (2010) cho hình chữ U ngược.
Khoảng tin cậy 95 \% được tính bằng phương pháp delta (delta method). Cả
bốn mẫu đều xác nhận chữ U ngược với điểm ngoặt trong phạm vi 39--46 \%
FSTS.

\textbf{Kiểm định ổn định hệ số giữa các đợt theo Paternoster (1998)
(M7):}

\begin{longtable}[]{@{}lllll@{}}
\toprule\noalign{}
So sánh & FSTSc z (p) & FSTSc² z (p) & \(\tilde{\mathrm{TCI}}\) z (p) &
\(\tilde{\mathrm{DAI}}\) z (p) \\
\midrule\noalign{}
\endhead
\bottomrule\noalign{}
\endlastfoot
2009 vs. 2015 & z=-0,61 (p=0,543) & z=0,96 (p=0,337) & z=1,29 (p=0,198)
& z=3,35 (p=0,001) \\
2009 vs. 2023 & z=0,09 (p=0,926) & z=0,14 (p=0,888) & z=1,42 (p=0,155) &
z=1,28 (p=0,201) \\
2015 vs. 2023 & z=0,66 (p=0,509) & z=-0,84 (p=0,400) & z=0,07 (p=0,947)
& z=-2,05 (p=0,040) \\
\end{longtable}

\emph{Ghi chú.} Kiểm định tính bằng nhau của hệ số giữa các đợt theo
Paternoster et al. (1998). Không bác bỏ chỉ ra tính ổn định của hệ số
giữa các đợt. So sánh FSTSc và FSTSc² ủng hộ tính không đồng nhất đặc
trưng theo đợt trong các tác động trực tiếp TCI và DAI nhưng không trong
các tham số độ cong I--P.

\begin{longtable}[]{@{}lllll@{}}
\toprule\noalign{}
Mẫu & N & Điểm ngoặt (FSTS thô) & 95 \% CI phương pháp delta &
Lind--Mehlum p \\
\midrule\noalign{}
\endhead
\bottomrule\noalign{}
\endlastfoot
Việt Nam 2009 & 989 & 46,2 \% & {[}37,4 \%, 55,1 \%{]} & .006 \\
Việt Nam 2015 & 956 & 39,3 \% & {[}30,3 \%, 48,4 \%{]} & .009 \\
Việt Nam 2023 & 1.013 & 41,6 \% & {[}31,7 \%, 51,5 \%{]} & .013 \\
Việt Nam gộp & 2.958 & 39,7 \% & {[}34,0 \%, 45,5 \%{]} & \textless{}
.001 \\
\end{longtable}

Các điểm ngoặt được rút ra từ đặc tả chữ U ngược M2
(\(\ln(\mathrm{LP})\) = β0 + β1 \(\mathrm{FSTS}^{c}\) + β2
\((\mathrm{FSTS}^{c})^2\) + biến kiểm soát + sector FE {[}+ wave FE
trong mẫu gộp{]}) và được biến đổi ngược về thang FSTS thô. Khoảng tin
cậy 95 \% sử dụng phương pháp delta trên phạm vi \(\mathrm{FSTS}^{c}\)
trong mẫu. Các giá trị p Lind--Mehlum theo kiểm định điểm cuối kiểu
Sasabuchi cho chữ U ngược trên phạm vi FSTS quan sát được (Lind and
Mehlum, 2010). Nguồn: Khảo sát Doanh nghiệp của Ngân hàng Thế giới; tính
toán của các tác giả.

\hypertarget{5-thux1ea3o-luux1eadn}{%
\subsection{5. Thảo luận}\label{5-thux1ea3o-luux1eadn}}

\hypertarget{51-diux1ec5n-giux1ea3i-lux1ea1i-nux103ng-lux1ef1c-sux1ed1-ux1edf-viux1ec7t-nam}{%
\subsubsection{5.1 Diễn giải lại năng lực số ở Việt
Nam}\label{51-diux1ec5n-giux1ea3i-lux1ea1i-nux103ng-lux1ef1c-sux1ed1-ux1edf-viux1ec7t-nam}}

Hàm ý trung tâm của các phát hiện là áp dụng công nghệ số nền tảng ở
Việt Nam không nên được diễn giải như một khoản phụ trội năng suất phổ
quát và ổn định theo thời gian. Mặc dù cả \(\tilde{\mathrm{TCI}}\) và
\(\tilde{\mathrm{DAI}}\) đều dương tính trung bình trong các đặc tả gộp,
vai trò thực nghiệm của chúng khác biệt đáng kể giữa các đợt và giữa các
chiến lược nhận diện (Vahlne, 2020; Stallkamp and Schotter, 2021). Điều
này có nghĩa là áp dụng công nghệ số dựa trên website không đơn giản là
một lợi thế nền không đổi; nó là một nguồn nhạy cảm với bối cảnh và phụ
thuộc giai đoạn của tính không đồng nhất hiệu quả.

Dấu âm của tương tác \(\mathrm{FSTS}^{c}\) × \(\tilde{\mathrm{DAI}}\)
vào năm 2023 (nơi tính liên quan đến năng suất của DAI suy giảm ở cường
độ xuất khẩu cao hơn) nhất quán với biên Tầng 1 của khái niệm. Cụ thể,
website đơn lẻ không thể quản lý mật độ giao dịch của các hoạt động
cường độ xuất khẩu cao mà không có hỗ trợ thanh toán điện tử và tích hợp
quy trình. Mẫu hình này tương phản với các bối cảnh nơi cùng khái niệm
bao gồm các chỉ báo kích hoạt giao dịch Tầng 2: ở đó, cường độ xuất khẩu
cao hơn có thể khuếch đại lợi nhuận năng suất từ áp dụng công nghệ số.

Cách diễn giải này giúp dung hoà sự cùng tồn tại của các tác động gộp
dương và các kết quả theo từng đợt không đồng đều. Mô hình gộp ghi nhận
xu hướng trung bình rằng năng lực mạnh hơn có liên hệ với hiệu quả tốt
hơn. Các mô hình theo từng đợt cho thấy xu hướng này không đồng đều mạnh
trong mọi giai đoạn. Giá trị của năng lực số do đó phụ thuộc vào vị trí
của doanh nghiệp trong vòng đời rộng hơn của quốc tế hoá và chuyển đổi.

\hypertarget{52-vuxec-sao-phuxe2n-biux1ec7t-giux1eefa-tci-vuxe0-dai-lux1ea1i-quan-trux1ecdng}{%
\subsubsection{5.2 Vì sao phân biệt giữa TCI và DAI lại quan
trọng}\label{52-vuxec-sao-phuxe2n-biux1ec7t-giux1eefa-tci-vuxe0-dai-lux1ea1i-quan-trux1ecdng}}

Các kết quả củng cố luận điểm lý thuyết về việc tách năng lực công nghệ
nước ngoài / tiêu chuẩn khỏi áp dụng công nghệ số nền tảng. Bằng chứng
PSM và IV trong mục 4.5 (Panel J và K) làm rõ hơn sự phân biệt (Karna et
al., 2016).

TCI vững chắc dưới cả đối sánh và công cụ hoá: ước lượng 2SLS của
\(\tilde{\mathrm{TCI}}\) là 1,64 (p \textless{} .001) dưới một công cụ
mạnh (giai đoạn một F = 22,1), và ATT đối sánh cho xử lý chứng nhận /
công nghệ nước ngoài là 0,61--0,64 (p \textless{} .001) trên một mẫu đối
sánh khoảng 640 doanh nghiệp.

Ngược lại, mối liên hệ trực tiếp DAI được OLS phát hiện được tái tạo
dưới PSM (ATT đối sánh = 0,30--0,32, p \textless{} .001) nhưng suy giảm
về không dưới 2SLS (β = 0,02, p = .94). Tóm lại, các kết quả gợi ý rằng
năng lực công nghệ nước ngoài / tiêu chuẩn cư xử như một kênh năng suất
vững hơn về nhận diện, trong khi hiện diện số dựa trên website cư xử như
một dấu hiệu nhạy cảm hơn với bối cảnh và chọn lựa của tính không đồng
nhất hiệu quả.

Phần mở rộng \(\mathrm{DAI}_{\text{rich}}\) ở mục 4.5 Panel B củng cố
cách diễn giải khái niệm thay vì làm suy yếu nó.
\(\tilde{\mathrm{DAI}}\) chính neo vào c22b (hiện diện website), đến năm
2023 gần với chỉ báo cơ sở Tầng 1. Phần mở rộng
\(\mathrm{DAI}_{\text{rich}}\) chỉ có vào năm 2023, kết hợp c22b với tỉ
trọng thanh toán điện tử (k33, k38). Phần mở rộng này tạo ra mẫu hình
điều tiết theo cùng hướng và có ý nghĩa cận ngưỡng tương tự
(\(\mathrm{FSTS}^{c}\) × DAI\_rich\_cont\_z = -0,93, M8 joint p = .099).
Dù áp dụng công nghệ số được đo bằng chỉ báo website Tầng 1 mỏng hay
bằng các chỉ báo giao dịch Tầng 2/3 sâu hơn, mẫu hình điều tiết năm 2023
đi theo cùng hướng. Bằng chứng cùng hướng này bảo vệ chống lại cách đọc
tính lỗi thời của biến đại diện cho mẫu hình DAI một cách riêng lẻ. Ngay
cả dưới thước đo phong phú hơn (khó bị bác bỏ như dấu hiệu thường lệ),
điều tiết trong năm 2023 vẫn nhất quán với phát hiện chính.

Sự khác biệt giữa TCI và DAI quan trọng vì hai khái niệm không cư xử
giống nhau qua các giai đoạn hoặc dưới các chiến lược nhận diện khác
nhau.

\(\tilde{\mathrm{TCI}}\) cư xử như một năng lực công nghệ nước ngoài /
tiêu chuẩn vững chắc với mối liên hệ năng suất sống sót qua cả PSM và
2SLS. \(\tilde{\mathrm{DAI}}\) cư xử như một dấu hiệu nhạy cảm với bối
cảnh với mối liên hệ năng suất được ghi nhận dưới so sánh đối sánh nhưng
không dưới biến thiên công cụ ngoại sinh. Điều này ủng hộ một cách tiếp
cận cẩn trọng hơn đối với số hoá trong nghiên cứu kinh doanh quốc tế,
một cách tiếp cận phân biệt giữa kích hoạt số cơ bản (dễ bị nhiễm bởi
chọn lựa trên các biến quan sát được) và chiều sâu công nghệ rộng hơn
(vững hơn về nhận diện), thay vì gộp chúng vào một nhãn duy nhất.

\hypertarget{53-uxfd-nghux129a-cux1ee7a-giai-ux111oux1ea1n-luxf5m-2015}{%
\subsubsection{5.3 Ý nghĩa của giai đoạn lõm
2015}\label{53-uxfd-nghux129a-cux1ee7a-giai-ux111oux1ea1n-luxf5m-2015}}

Trước khi viện dẫn câu chuyện hạ tầng vĩ mô, cần neo kết quả không có ý
nghĩa năm 2015 vào cấu trúc vi mô của bản thân mẫu. Như báo cáo trong
Bảng 1, tỉ lệ doanh nghiệp xuất khẩu giảm từ 28,4 \% năm 2009 xuống 20,7
\% năm 2015. Đây là sự rút lui có chọn lọc, để lại nhóm tồn tại tập
trung không cân đối ở các hoạt động lắp ráp thâm dụng lao động và phụ
thuộc OEM. Đối với nhà thầu phụ dệt may hoặc điện tử đang chờ đơn hàng
từ nguyên đơn nước ngoài, website doanh nghiệp cơ bản vận hành như tờ
bướm số chứ không phải công cụ phối hợp nâng cao năng suất. Nó phát tín
hiệu hiện diện nhưng không tối ưu hoá dòng chuỗi cung ứng, không kích
hoạt tích hợp dữ liệu ở cấp giao dịch, hay thay thế các giao diện số do
người mua thượng nguồn kiểm soát. Cơ chế thành phần nhóm này (không chỉ
khoảng cách hạ tầng vĩ mô) giúp giải thích vì sao hệ số DAI dựa trên
website nén về không vào năm 2015, ngay cả khi độ cong chữ U ngược vẫn
còn ý nghĩa thống kê. Kênh tính lỗi thời của biến đại diện củng cố cách
đọc này. Đến năm 2015, sở hữu website đã lan toả đủ để chỉ báo c22b
không còn có thể tách biệt các doanh nghiệp xuất khẩu có năng lực số
khỏi sự hiện diện web thường lệ, làm suy giảm thêm mối liên hệ đo lường
được với năng suất. Do đó, các chỉ báo vĩ mô dưới đây được xem như bối
cảnh thể chế khiến cách đọc theo từng đợt trở nên khả tín, chứ không
phải như biến thiên nhận diện.

Một cách diễn giải thuần tuý về phụ thuộc giai đoạn không phải là cách
đọc duy nhất của mẫu hình DAI theo từng đợt. Một giải thích thay thế, và
ở một mức độ là bổ sung, là tính lỗi thời của biến đại diện: chỉ báo
c22b về sở hữu website đo lường thứ gì đó khác nhau trong các năm 2009,
2015 và 2023 ngay cả khi văn bản câu hỏi giống nhau giữa các đợt.

Năm 2009, sở hữu website doanh nghiệp là giao diện thị trường tương đối
khác biệt, phát tín hiệu về tính hợp pháp số cơ bản trong giao dịch quốc
tế. Đến năm 2023, website gần với dấu hiệu thường lệ của hiện diện số cơ
bản, được chia sẻ giữa các doanh nghiệp xuất khẩu và không xuất khẩu.
Dưới cách đọc này, sự suy giảm mối liên hệ trực tiếp DAI giữa các đợt
phần nào phản ánh nội dung thông tin co lại của c22b, chứ không phải
thay đổi cấu trúc trong mức sinh lời năng suất đối với áp dụng công nghệ
số nền tảng. Nghiên cứu không xem tính lỗi thời của biến đại diện và phụ
thuộc giai đoạn là loại trừ lẫn nhau: cả hai có thể đang vận hành đồng
thời. Nghiên cứu nhấn mạnh cách đọc tính lỗi thời của biến đại diện vì
nó không thể kiểm định trong công cụ WBES có thể so sánh giữa các đợt.
Công trình tương lai sử dụng các đợt với các chỉ báo số Tầng 2--3 phong
phú hơn sẽ cần thiết để tách rời hai kênh.

Mẫu hình 2015 được diễn giải tốt nhất không phải như bất thường, mà như
sự nén đặc trưng theo đợt của kênh áp dụng công nghệ số nền tảng dưới
điều kiện hạ tầng chuyển đổi. Mẫu này cho thấy ngay cả khi cấu trúc I--P
phi tuyến trở nên rõ ràng hơn, mức sinh lời trực tiếp từ hiện diện số
dựa trên website có thể nén về không. Cách đọc này khả tín về mặt thể
chế, nhưng không nên được xem như giải thích chính thức được nhận diện
về mẫu hình hệ số. Doanh nghiệp có thể đi qua giai đoạn chuyển đổi trong
đó mở rộng xuất khẩu vẫn quan trọng nhưng đóng góp năng suất của áp dụng
công nghệ số cơ bản trở nên khó nhận thấy hoặc phát hiện trong mẫu hơn.
Sự nén năm 2015 được đối xử thẳng thắn như mối liên hệ đặc trưng theo
đợt nhất quán với phụ thuộc giai đoạn, chứ không phải dịch chuyển cấu
trúc được nhận diện đầy đủ giữa các đợt. Như báo cáo trong mục 4.5 Panel
I, kiểm định tương tác chính thức gộp đợt × biến trọng tâm chỉ phát hiện
các dịch chuyển trực tiếp DAI là có thể phân biệt giữa các đợt. Các khác
biệt độ cong FSTS và điều tiết FSTS × DAI giữa các đợt không thể tách
biệt thống kê.

Bối cảnh vĩ mô của giai đoạn nhất quán với vùng trũng hạ tầng số vào năm
2015 so với các neo 2009 và 2023. Các chỉ báo thứ cấp công khai ủng hộ
cách đọc này. Các Chỉ số Phát triển Thế giới (WDI) của Ngân hàng Thế
giới ghi nhận tỉ lệ cá nhân sử dụng Internet của Việt Nam ở mức khoảng
26 \% năm 2009, 45 \% năm 2015 và hơn 78 \% vào năm 2023. Thuê bao băng
rộng cố định trên 100 người của ITU tăng từ khoảng 3 năm 2009 lên 8 năm
2015 và hơn 20 vào năm 2023.

Chương trình Chuyển đổi Số Quốc gia (NDTP) của Việt Nam được ban hành
năm 2020. Các nền tảng thanh toán điện tử xuyên biên giới (VNPAY, MoMo,
ZaloPay) và các thị trường xuất khẩu B2B (Alibaba.com, Amazon Global
Selling) đạt quy mô tại Việt Nam chỉ sau 2018--2019. Đợt 2015 do đó quan
sát các doanh nghiệp xuất khẩu Việt Nam trong giai đoạn hạ tầng chuyển
đổi: website chưa thể cắm vào hệ sinh thái số hỗ trợ giao dịch. Đợt 2023
quan sát họ sau khi khung sườn hậu NDTP đã trưởng thành. Các chỉ báo vĩ
mô này không được xem là biến thiên nhận diện và không được đưa vào hồi
quy. Tuy vậy, chúng làm cho cách đọc theo từng đợt trở nên khả tín về
thể chế, chứ không thuần tuý hậu kiểm.

Thay vì xem đợt này như một bất thường, hữu ích hơn khi diễn giải nó như
bằng chứng về tính không đồng nhất đặc trưng theo đợt nhất quán với phụ
thuộc giai đoạn. Giai đoạn lõm chứng minh vì sao chỉ các mức trung bình
gộp là không đủ. Không có phân tích theo từng đợt, người ta sẽ bỏ qua
khả năng các khoản sinh lời năng lực nén hoặc mờ đi tạm thời trước khi
tái xuất hiện trong giai đoạn sau, ngay cả khi các tham số độ cong của
bản thân mối quan hệ I--P vẫn không thể phân biệt thống kê giữa các đợt.

\hypertarget{54-huxe0m-uxfd-thux1ef1c-tiux1ec5n}{%
\subsubsection{5.4 Hàm ý thực
tiễn}\label{54-huxe0m-uxfd-thux1ef1c-tiux1ec5n}}

Hàm ý quản trị không phải là doanh nghiệp nên đầu tư ít hơn vào số hoá.
Đúng hơn, đó là số hoá nên được liên kết với mức độ sẵn sàng của tổ chức
và quy mô quốc tế. Các công cụ số nền tảng có thể tạo ra giá trị, nhưng
mức sinh lời phụ thuộc vào việc các công cụ đó có được nhúng vào các
thông lệ hỗ trợ phối hợp xuất khẩu và hiệu quả hay không.

Các nhà quản lý cũng nên tránh nhầm lẫn áp dụng công nghệ số cơ bản với
năng lực công nghệ sâu hơn.

Một doanh nghiệp có thể áp dụng các công cụ số mà không có năng lực học
hỏi và nâng cấp cần thiết để duy trì lợi ích hiệu quả theo thời gian.

Ngược lại, năng lực nội bộ mạnh hơn có thể cho phép doanh nghiệp chuyển
hoá việc áp dụng công nghệ số tương đối cơ bản thành các kết quả có ý
nghĩa hơn. Sự phân biệt này đặc biệt quan trọng trong các nền kinh tế
chuyển đổi, nơi doanh nghiệp đi qua các giai đoạn phát triển năng lực
không đồng đều.

\hypertarget{56-huxe0m-uxfd-chuxednh-suxe1ch}{%
\subsubsection{5.6 Hàm ý chính
sách}\label{56-huxe0m-uxfd-chuxednh-suxe1ch}}

Chúng tôi đóng khung cách đọc chính sách như những cân nhắc tạm thời chứ
không phải các quy định chỉ thị. Tính chất liên hệ của bằng chứng, độ
rộng của tính không đồng nhất theo từng đợt, và phạm vi một nền kinh tế
đơn lẻ đều cân nhắc chống lại việc chuyển các phát hiện thành mục tiêu
chính sách doanh nghiệp cụ thể. Với những lưu ý đó được giữ rõ ràng, ba
cân nhắc tiếp theo cho thiết kế chính sách thương mại và kinh tế số của
Việt Nam.

Thứ nhất, các công cụ thúc đẩy xuất khẩu được thiết kế quanh một khoản
phụ trội quốc tế hoá dương đồng nhất sẽ vượt mức trong các giai đoạn
chuyển đổi như đợt 2015, khi các khoản sinh lời năng lực bị nén ngay cả
khi cường độ xuất khẩu vẫn liên quan đến năng suất. Hướng hỗ trợ xuất
khẩu vào các doanh nghiệp nằm trong vùng chi phí nhập cuộc, thay vì các
doanh nghiệp đã hoạt động ở cường độ xuất khẩu cao, nhất quán với độ
cong được ghi nhận trong mẫu gộp và mẫu đợt sau. Thứ hai, các chương
trình kinh tế số xem áp dụng nền tảng (website, thanh toán điện tử cơ
bản) như một đòn bẩy chính sách đủ có khả năng bị suy giảm bởi độ trễ
triển khai và bởi tính có điều kiện của kênh số ở cường độ xuất khẩu cao
hơn.

Các chương trình kết hợp áp dụng công nghệ số Tầng 1--2 với nâng cấp
năng lực sâu hơn, chứng nhận chất lượng, đầu tư năng lực hấp thụ, các
thông lệ tổ chức cho phối hợp xuyên biên giới, sẽ nhất quán hơn với mẫu
hình nổi lên vào năm 2023.

Thứ ba, cách đọc vòng đời gợi ý rằng các cửa sổ đánh giá chính sách có ý
nghĩa quan trọng. Một chương trình chuyển đổi số chỉ được đánh giá dựa
trên đường cơ sở chuyển đổi kiểu 2015 sẽ đánh giá thấp đóng góp năng
suất dài hạn của nó; một chương trình được đánh giá dựa trên đường cơ sở
kiểu 2009 hoặc 2023 sẽ đánh giá cao so với giai đoạn ở giữa. Đánh giá
chính sách xem xét nghiêm túc cấu trúc vòng đời sẽ ghép đo lường kết quả
cửa sổ ngắn với đo lường bền vững về môi trường năng lực và hạ tầng nơi
doanh nghiệp hoạt động.

\hypertarget{6-hux1ea1n-chux1ebf-vuxe0-hux1b0ux1edbng-nghiuxean-cux1ee9u-tiux1ebfp-theo}{%
\subsection{6. Hạn chế và hướng nghiên cứu tiếp
theo}\label{6-hux1ea1n-chux1ebf-vuxe0-hux1b0ux1edbng-nghiuxean-cux1ee9u-tiux1ebfp-theo}}

Các phát hiện nên được đọc trong bối cảnh năm hạn chế. Thứ nhất và cơ
bản nhất, dữ liệu vi mô WBES là cắt ngang lặp lại chứ không phải một
bảng dữ liệu doanh nghiệp thực sự. Thay đổi nội tại doanh nghiệp theo
thời gian không thể được nhận diện, và tính không đồng nhất không quan
sát được bất biến theo thời gian không thể được loại trừ ở cấp doanh
nghiệp. Ngôn ngữ liên hệ được sử dụng xuyên suốt bài báo phản ánh ràng
buộc này và không nên được nới lỏng trong bất kỳ diễn giải nào của người
đọc về các kết quả.

Thứ hai, \(\tilde{\mathrm{DAI}}\) chỉ ghi nhận một lớp áp dụng công nghệ
số nền tảng kiểu Tầng 1 tập trung vào hiện diện website, chứ không phải
năng lực tổ chức được tích hợp số. Phần mở rộng
\(\mathrm{DAI}_{\text{rich}}\) cho đợt 2023 cho thấy mẫu hình có điều
kiện tồn tại khi các chỉ báo phong phú hơn có sẵn. Tuy nhiên, một bảng
dữ liệu các doanh nghiệp được tích hợp số với đo lường sâu hơn sẽ thắt
chặt khái niệm hơn và cho phép kiểm định sắc nét hơn về việc liệu tích
hợp số sâu hơn có thể hiện kênh năng suất ổn định hơn hay không.

Thứ ba, phân tích được thực hiện trên một nền kinh tế chuyển đổi đơn lẻ.
Việt Nam có giá trị thông tin chính vì môi trường thể chế và số của nó
dịch chuyển đáng kể trên cửa sổ quan sát 2009--2023. Tuy nhiên, mẫu hình
vòng đời được ghi nhận ở đây có thể không khái quát hoá mà không có sửa
đổi cho các nền kinh tế có hạ tầng số hoặc thành phần xuất khẩu theo quỹ
đạo khác.

Thứ tư, bằng chứng giữa các đợt không đồng đều về sức mạnh thống kê.

Các kiểm định z Paternoster et al. (1998) xác nhận rằng sự giảm
\(\tilde{\mathrm{DAI}}\) giữa 2009 và 2015 (z = 3,353, p \textless{}
.001) và sự phục hồi giữa 2015 và 2023 (z = -2,051, p = .040) có thể
phân biệt thống kê. Tuy vậy, hầu hết các khác biệt hệ số giữa các đợt
khác nằm ở độ lớn cận ngưỡng hoặc không có ý nghĩa. Diễn giải vòng đời
do đó dựa chủ yếu vào tính nhất quán định hướng của các ước lượng theo
từng đợt và sự tập trung của tín hiệu số vào năm 2023, chứ không dựa vào
các khác biệt hệ số cặp đôi có ý nghĩa đồng đều qua tất cả các số hạng
trọng tâm. Công trình tương lai có thể khai thác thời điểm chính sách
(ví dụ, Chương trình Chuyển đổi Số Quốc gia của Việt Nam khởi động năm
2020) để nhận diện sắc nét hơn kênh số thông qua thiết kế đánh giá chính
sách thay vì so sánh giữa các đợt.

Thứ năm, các hiệu ứng cố định ngành trong các mô hình chính được cố ý
rộng để giữ đặc tả giữa các đợt có thể so sánh. Mục 4.5 Panel G giải
quyết một phần mở rộng tự nhiên bằng cách tái ước lượng các đặc tả
M2/M7/M8 gộp trên các nhóm phụ sản xuất (sector1 ∈ \{1, 2, 3\}; N =
1.854) so với phi sản xuất (sector1 ∈ \{4, 5, 6, 7\}; N = 1.104). Kết
quả cho thấy kênh áp dụng công nghệ số vận hành chủ yếu ở sản xuất,
trong khi kênh năng lực công nghệ hiện diện trong cả hai. Các cơ chế cấp
ngành tinh hơn (ví dụ, tương phản các dịch vụ thâm dụng số với các dịch
vụ truyền thống, hoặc các phân ngành sản xuất định hướng xuất khẩu với
các phân ngành định hướng nội địa) sẽ được hưởng lợi từ thiết kế nghiên
cứu khai thác phân loại ngành phong phú hơn so với các mã WBES một chữ
số.

Thứ sáu, phân tích được giới hạn ở Việt Nam như nền kinh tế chuyển đổi
đơn lẻ ở giai đoạn cụ thể của phát triển số và thể chế. Mẫu hình điều
tiết DAI được ghi nhận ở đây (âm ở cường độ xuất khẩu cao dưới áp dụng
chỉ riêng Tầng 1) cần được so sánh với hai loại bối cảnh khác: các thị
trường mới nổi khác ở các giai đoạn thể chế khác nhau, và các nền kinh
tế tiên tiến về số nơi tương tác tương tự được dự đoán là dương. So sánh
như vậy đòi hỏi các thiết kế giữa các bối cảnh thay đổi một cách hệ
thống chi phí giao dịch thể chế và chất lượng hạ tầng số. Các thiết kế
như vậy cũng sẽ cho phép kiểm định trực tiếp hơn về khung phụ thuộc bối
cảnh ở Mục 2.1.

\hypertarget{7-kux1ebft-luux1eadn}{%
\subsection{7. Kết luận}\label{7-kux1ebft-luux1eadn}}

Nghiên cứu này xem xét lại mối quan hệ quốc tế hoá--hiệu quả ở Việt Nam
bằng cách phân biệt năng lực công nghệ nước ngoài / tiêu chuẩn khỏi áp
dụng công nghệ số nền tảng và bằng cách so sánh bằng chứng gộp với bằng
chứng theo từng đợt.

Các phát hiện cho thấy mối quan hệ quốc tế hoá--hiệu quả là không đơn
điệu trong toàn mẫu, nhưng mẫu hình này được chi phối chủ yếu bởi cấu
trúc tham gia-và-cường độ chứ không phải bởi độ cong mạnh trong nhóm
doanh nghiệp xuất khẩu đơn lẻ.

Các kết quả cũng cho thấy hai miền năng lực không nên được xem là có thể
hoán đổi. Năng lực công nghệ tương đối ổn định và vững về nhận diện qua
các đặc tả, trong khi áp dụng công nghệ số nền tảng nhạy cảm theo đợt
hơn và chỉ thể hiện điều tiết có thể phát hiện trong mẫu vào năm 2023.
Hàm ý rộng hơn là: tính liên quan đến năng suất của áp dụng công nghệ số
cơ bản trong nền kinh tế chuyển đổi là thực sự nhưng không đồng đều. Các
mức trung bình gộp có thể che giấu tính không đồng nhất theo thời gian
quan trọng về việc khi nào và bằng cách nào áp dụng công nghệ số có ý
nghĩa.

\hypertarget{xung-ux111ux1ed9t-lux1ee3i-uxedch}{%
\subsection{Xung đột lợi ích}\label{xung-ux111ux1ed9t-lux1ee3i-uxedch}}

Các tác giả tuyên bố không có xung đột lợi ích.

\hypertarget{tuxe0i-trux1ee3}{%
\subsection{Tài trợ}\label{tuxe0i-trux1ee3}}

Nghiên cứu này không nhận được tài trợ cụ thể nào từ bất kỳ cơ quan tài
trợ nào trong khu vực công, thương mại hoặc phi lợi nhuận.

\hypertarget{tuyuxean-bux1ed1-vux1ec1-sux1eb5n-cuxf3-dux1eef-liux1ec7u}{%
\subsection{Tuyên bố về sẵn có dữ
liệu}\label{tuyuxean-bux1ed1-vux1ec1-sux1eb5n-cuxf3-dux1eef-liux1ec7u}}

Dữ liệu hỗ trợ các phát hiện của nghiên cứu này là từ Khảo sát Doanh
nghiệp của Ngân hàng Thế giới và có sẵn từ cổng thông tin Khảo sát Doanh
nghiệp của Ngân hàng Thế giới, tuỳ thuộc vào việc đăng ký và tuân thủ
Giao thức Tiếp cận Dữ liệu Khảo sát Doanh nghiệp của Ngân hàng Thế giới.
Vì giao thức hạn chế việc tái phân phối các tệp .dta gốc cho bên thứ ba,
các tác giả không tái phân phối các tệp khảo sát thô. Tài liệu tái lặp,
bao gồm chi tiết xây dựng biến, đặc tả mô hình và đầu ra tính toán, có
sẵn từ các tác giả khi có yêu cầu hợp lý và có thể được chia sẻ trong
phạm vi cho phép của các điều khoản tiếp cận dữ liệu.

\hypertarget{sux1eed-dux1ee5ng-ai-tux1ea1o-sinh-trong-quuxe1-truxecnh-viux1ebft}{%
\subsection{Sử dụng AI tạo sinh trong quá trình
viết}\label{sux1eed-dux1ee5ng-ai-tux1ea1o-sinh-trong-quuxe1-truxecnh-viux1ebft}}

Các công cụ AI tạo sinh đã được sử dụng trong quá trình chuẩn bị bản
thảo để hỗ trợ chỉnh sửa ngôn ngữ, gợi ý cấu trúc và lắp ráp tài liệu
gói tái lặp.

Toàn bộ việc đóng khung khái niệm, phát triển giả thuyết, phân tích thực
nghiệm, diễn giải kết quả và hoàn thiện ngôn từ đều do các tác giả con
người thực hiện, những người chịu hoàn toàn trách nhiệm về nội dung của
công trình xuất bản.

\hypertarget{lux1eddi-cux1ea3m-ux1a1n}{%
\subsection{Lời cảm ơn}\label{lux1eddi-cux1ea3m-ux1a1n}}

Nghiên cứu này sử dụng dữ liệu từ Khảo sát Doanh nghiệp của Ngân hàng
Thế giới. Chúng tôi cảm ơn Đơn vị Phân tích Doanh nghiệp của Nhóm Chỉ số
Toàn cầu Kinh tế Phát triển thuộc Ngân hàng Thế giới đã cung cấp quyền
tiếp cận dữ liệu.

Người thu thập dữ liệu gốc, nhà phân phối được uỷ quyền và cơ quan tài
trợ liên quan không chịu trách nhiệm về việc sử dụng dữ liệu hoặc các
diễn giải dựa trên việc sử dụng đó. Các phát hiện, diễn giải và kết luận
được trình bày trong bài báo này là của các tác giả và không nhất thiết
đại diện cho quan điểm của Nhóm Ngân hàng Thế giới, các Giám đốc Điều
hành hoặc các chính phủ mà họ đại diện. Các tác giả không nhận được tài
trợ cụ thể nào từ bất kỳ cơ quan tài trợ nào trong khu vực công, thương
mại hoặc phi lợi nhuận.

Gói tái lặp đi kèm bản thảo này được phát triển để hỗ trợ tính minh bạch
về phương pháp luận và khả năng tái tạo tính toán.

\hypertarget{tuxe0i-liux1ec7u-bux1ed5-sung}{%
\subsection{Tài liệu bổ sung}\label{tuxe0i-liux1ec7u-bux1ed5-sung}}

Tất cả các hình và bảng trong bản thảo đều được nhúng trực tiếp. Các đầu
ra hệ số nền cũng được phân phối dưới dạng các tệp CSV thuần trong gói
tái lặp để phân tích lại tiếp theo:

tables/table\_1\_descriptives.csv (Bảng 1),
tables/coefs\_main\_models.csv (bảng hệ số định dạng dài M0--M8),
tables/joint\_tests\_main\_models.csv (các kiểm định F liên hợp H2 và
P1), tables/table\_lind\_mehlum.csv (nguồn Bảng LM),
tables/table\_3\_robustness.csv (các bảng A--D + G của Bảng 3),
tables/selection\_checks.csv (Heckman / hàm kiểm soát), và
tables/table\_paternoster.csv (các kiểm định z giữa các đợt). Gói tái
lặp và hướng dẫn: xem p3\_vietnam/README.md.

\begin{center}\rule{0.5\linewidth}{0.5pt}\end{center}

\hypertarget{references}{%
\subsection{References}\label{references}}

Aguinis, H., Hill, N. S., and Bailey, J. R. (2021). Best practices in
data collection and preparation: Recommendations for reviewers, editors,
and authors. Organizational Research Methods, 24(4):678--693.

Antonakis, J., Bendahan, S., Jacquart, P., and Lalive, R. (2010). On
making causal claims: A review and recommendations. The Leadership
Quarterly, 21(6):1086--1120.

Author Citation (2026 --- JFAR). Unveiling the impact of Chinese
manufacturing SMEs\textquotesingle{} internationalization on
performance. \emph{Journal of Finance \& Accounting Research, 02}(39),
287--291. {[}Author details withheld for blind review.{]}

Banalieva, E. R. and Dhanaraj, C. (2019). Internalization theory for the
digital economy. Journal of International Business Studies,
50(8):1372--1387.

Bernard, A. B., Jensen, J. B., Redding, S. J., and Schott, P. K. (2007).
Firms in international trade. Journal of Economic Perspectives,
21(3):105--130.

Bharadwaj, A., El Sawy, O. A., Pavlou, P. A., and Venkatraman, N.
(2013). Digital business strategy: Toward a next generation of insights.
MIS Quarterly, 37(2):471--482.

Brouthers, K. D., Geisser, K. D., and Rothlauf, F. (2016). Explaining
the internationalization of ibusiness firms. Journal of International
Business Studies, 47(5):513--534.

Brynjolfsson, E., Rock, D., and Syverson, C. (2021). The productivity
J-curve: How intangibles complement general purpose technologies.
American Economic Journal: Macroeconomics, 13(1):333--372.

Certo, S. T., Busenbark, J. R., Woo, H.-s., and Semadeni, M. (2016). and
Heckman models in strategic management research. Sample selection bias
Strategic Management Journal, 37(13):2639--2657.

Cohen, W. M. and Levinthal, D. A. (1990). Absorptive capacity: A new
perspective on learning and innovation. Administrative Science
Quarterly, 35(1):128--152.

Contractor, F. J. (2007). Is international business good for companies?
The evolutionary or multi-stage theory of internationalization vs. the
transaction cost perspective. Management International Review,
47(3):453--475.

Coviello, N., Kano, L., and Liesch, P. W. (2017). Adapting the Uppsala
model to a modern world: Macro-context and microfoundations. Journal of
International Business Studies, 48(9):1151-- 1164. Cuervo-Cazurra, A.
and Genc, M. (2008). Transforming disadvantages into advantages:
Developing-country MNEs in the least developed countries. Journal of
International Business Studies, 39(6):957--979.

Dang, T. M. D., Le, V., \& Nguyen, T. T. (2024). Export persistence and
productivity in Vietnam: Who learns? \emph{Journal of Economic
Integration}, \emph{39}(2), 234--262.

Eisenhardt, K. M. and Martin, J. A. (2000). Dynamic capabilities: What
are they? Strategic Management Journal, 21(10--11):1105--1121.

Goldfarb, A. and Tucker, C. (2019). Digital economics. Journal of
Economic Literature, 57(1):3-- 43.

Haans, R. F. J., Pieters, C., and He, Z.-L. (2016). Thinking about U:
Theorizing and testing Uand inverted-U-shaped relationships in strategy
research. Strategic Management Journal, 37(7):1177--1195.

Hanelt, A., Bohnsack, R., Marz, D., and Antunes Marante, C. (2021). A
systematic review of the literature on digital transformation: Insights
and implications for strategy and organizational change. Journal of
Management Studies, 58(5):1159--1197.

Heckman, J. J. (1979). Sample selection bias as a specification error.
Econometrica, 47(1):153-- 161.

Hennart, J.-F. (2007). The theoretical rationale for a
multinationality--performance relationship. Management International
Review, 47(3):423--452.

Hennart, J.-F. (2011). A theoretical assessment of the empirical
literature on the impact of multinationality on performance. Global
Strategy Journal, 1(1--2):135--151.

Karna, A., Richter, A., and Riesenkampff, E. (2016). Revisiting the role
of the environment in the capabilities--financial performance
relationship: A meta-analysis. Strategic Management Journal,
37(6):1154--1173.

Knight, G. A. and Liesch, P. W. (2016). Internationalization: From
incremental to born global. Journal of World Business, 51(1):93--102.

Lall, S. (1992). Technological capabilities and industrialization. World
Development, 20(2):165--186.

Lind, J. T. and Mehlum, H. (2010). With or without U? The appropriate
test for a U-shaped relationship. Oxford Bulletin of Economics and
Statistics, 72(1):109--118.

Lu, J. W. and Beamish, P. W. (2004). International diversification and
firm performance: The S-curve hypothesis. Academy of Management Journal,
47(4):598--609.

Marano, V., Arregle, J.-L., Hitt, M. A., Spadafora, E., and van Essen,
M. (2016). Home country institutions and the
internationalization--performance relationship: A meta-analytic review.
Journal of Management, 42(5):1075--1110.

Nambisan, S., Wright, M., and Feldman, M. (2019). The digital
transformation of innovation and entrepreneurship: Progress, challenges
and key themes. Research Policy, 48(8):103773.

Ngo, V. D., Nguyen, T. T., \& Hoang, T. T. (2024). Export-specific
investments, competitive advantage, and performance in Vietnamese SMEs:
The moderating role of domestic market conditions. \emph{Journal of
Business Research}, \emph{176}, 114577.

Nguyen, D. P. K., Phan, T. M., \& Le, T. T. H. (2025). The simultaneous
relationship between labor productivity and exports: The case of small
and medium-sized enterprises in Vietnam. \emph{Ho Chi Minh City Open
University Journal of Science: Economics and Business Administration},
\emph{20}(1), 1--18.

Oster, E. (2019). Unobservable selection and coefficient stability:
Theory and evidence. Journal of Business \& Economic Statistics,
37(2):187--204.

Paternoster, R., Brame, R., Mazerolle, P., and Piquero, A. (1998). Using
the correct statistical test for the equality of regression
coefficients. Criminology, 36(4):859--866.

Petricevic, O. and Teece, D. J. (2019). The structural reshaping of
globalization: Implications for strategic sectors, profiting from
innovation, and the multinational enterprise. Journal of International
Business Studies, 50(9):1487--1512.

Pisani, N., Garcia-Bernardo, J., and Heemskerk, E. (2020). Does it pay
to be a multinational? A large-sample, cross-national replication
assessing the multinationality--performance relationship. Strategic
Management Journal, 41(1):152--172.

Shaver, J. M. (2020). Causal identification through a cumulative body of
research in the study of strategy and organizations. Journal of
Management, 46(7):1244--1256.

Stallkamp, M. and Schotter, A. P. J. (2021). Platforms without borders?
The international strategies of digital platform firms. Global Strategy
Journal, 11(1):58--80.

Strange, R. and Zucchella, A. (2017). Industry 4.0, global value chains
and international business. Multinational Business Review,
25(3):174--184.

Teece, D. J. (2007). Explicating dynamic capabilities: The nature and
microfoundations of (sustainable) enterprise performance. Strategic
Management Journal, 28(13):1319--1350.

Tran, T. T., Le, T. H., \& Pham, V. K. (2025). The impact of
innovation-driven digital transformation on export performance of SMEs.
\emph{Emerging Science Journal}, \emph{9}(2), 451--467.

Vahlne, J.-E. (2020). Development of the Uppsala model of
internationalization process: From internationalization to evolution.
Global Strategy Journal, 10(2):239--250.

Vahlne, J.-E. and Johanson, J. (2017). From internationalization to
evolution: The Uppsala model at 40 years. Journal of International
Business Studies, 48(9):1087--1102.

Verhoef, P. C., Broekhuizen, T., Bart, Y., Bhattacharya, A., Dong, J.
Q., Fabian, N., and Haenlein, M. (2021). Digital transformation: A
multidisciplinary reflection and research agenda. Journal of Business
Research, 122:889--901.

Vernon, R. (1979). The product cycle hypothesis in a new international
environment. Oxford Bulletin of Economics and Statistics,
41(4):255--267.

Vial, G. (2019). Understanding digital transformation: A review and a
research agenda. Journal of Strategic Information Systems,
28(2):118--144.

Wagner, J. (2007). Exports and productivity: A survey of the evidence
from firm-level data. The World Economy, 30(1):60--82.

Williamson, O. E. (1985). \emph{The economic institutions of
capitalism}. Free Press.

Wolfolds, S. E. and Siegel, J. (2019). Misaccounting for endogeneity:
The peril of relying on the Heckman two-step method without a valid
instrument. Strategic Management Journal, 40(3):432--462.

Wooldridge, J. M. (2010). Econometric analysis of cross section and
panel data. MIT Press, 2nd edition.

World Bank (2010). Vietnam Enterprise Survey 2009: Data file. Enterprise
Surveys, World Bank Group.

World Bank (2016). Vietnam Enterprise Survey 2015: Data file. Enterprise
Surveys, World Bank Group.

World Bank (2024). Vietnam Enterprise Survey 2023: Data file. Enterprise
Surveys, World Bank Group.

Wright, M., Filatotchev, I., Hoskisson, R. E., and Peng, M. W. (2005).
Strategy research in emerging economies: Challenging the conventional
wisdom. Journal of Management Studies, 42(1):1--33.

Wu, F., Fan, D., and Chen, L. (2022). Untangling the
internationalization--performance relationship in emerging-economy
multinationals: A meta-analytic review. Management International Review,
62(3):315--354.

Wu, J., Wang, C., Hong, J., Piperopoulos, P., and Zhuo, S. (2016).
Internationalization and innovation performance of emerging market
enterprises: The role of host-country institutional development. Journal
of World Business, 51(2):251--263.

\end{document}
